%%%%%%%% ICML 2025 EXAMPLE LATEX SUBMISSION FILE %%%%%%%%%%%%%%%%%

\documentclass{article}

% Recommended, but optional, packages for figures and better typesetting:
\usepackage{microtype}
\usepackage{graphicx}
\usepackage{subfigure}
\usepackage{booktabs} % for professional tables
% \usepackage{subcaption}
\usepackage{multirow}

% hyperref makes hyperlinks in the resulting PDF.
% If your build breaks (sometimes temporarily if a hyperlink spans a page)
% please comment out the following usepackage line and replace
% \usepackage{icml2025} with \usepackage[nohyperref]{icml2025} above.
\usepackage{hyperref}


% Attempt to make hyperref and algorithmic work together better:
\newcommand{\theHalgorithm}{\arabic{algorithm}}

% Use the following line for the initial blind version submitted for review:
\usepackage{icml2025}

% If accepted, instead use the following line for the camera-ready submission:
% \usepackage[accepted]{icml2025}

% For theorems and such
\usepackage{amsmath}
\usepackage{amssymb}
\usepackage{mathtools}
\usepackage{amsthm}

% if you use cleveref..
\usepackage[capitalize,noabbrev]{cleveref}

%%%%%%%%%%%%%%%%%%%%%%%%%%%%%%%%
% THEOREMS
%%%%%%%%%%%%%%%%%%%%%%%%%%%%%%%%
\theoremstyle{plain}
\newtheorem{theorem}{Theorem}[section]
\newtheorem{proposition}[theorem]{Proposition}
\newtheorem{lemma}[theorem]{Lemma}
\newtheorem{corollary}[theorem]{Corollary}
\theoremstyle{definition}
\newtheorem{definition}[theorem]{Definition}
\newtheorem{assumption}[theorem]{Assumption}
\theoremstyle{remark}
\newtheorem{remark}[theorem]{Remark}

% Todonotes is useful during development; simply uncomment the next line
%    and comment out the line below the next line to turn off comments
%\usepackage[disable,textsize=tiny]{todonotes}
\usepackage[textsize=tiny]{todonotes}


% The \icmltitle you define below is probably too long as a header.
% Therefore, a short form for the running title is supplied here:
\icmltitlerunning{Flat Minima in the Adapter Subspace Matters for Continual Learning}

\begin{document}

\twocolumn[
\icmltitle{Flat Minima in the Adapter Subspace Matters for Continual Learning: A PAC-Bayesian Perspective}

% It is OKAY to include author information, even for blind
% submissions: the style file will automatically remove it for you
% unless you've provided the [accepted] option to the icml2025
% package.

% List of affiliations: The first argument should be a (short)
% identifier you will use later to specify author affiliations
% Academic affiliations should list Department, University, City, Region, Country
% Industry affiliations should list Company, City, Region, Country

% You can specify symbols, otherwise they are numbered in order.
% Ideally, you should not use this facility. Affiliations will be numbered
% in order of appearance and this is the preferred way.
% \icmlsetsymbol{equal}{*}

\begin{icmlauthorlist}
% % \icmlauthor{Firstname1 Lastname1}{equal,yyy}
% % \icmlauthor{Firstname2 Lastname2}{equal,yyy,comp}
% % \icmlauthor{Firstname3 Lastname3}{comp}
% % \icmlauthor{Firstname4 Lastname4}{sch}
% % \icmlauthor{Firstname5 Lastname5}{yyy}
% % \icmlauthor{Firstname6 Lastname6}{sch,yyy,comp}
% % \icmlauthor{Firstname7 Lastname7}{comp}
% % %\icmlauthor{}{sch}
% % \icmlauthor{Firstname8 Lastname8}{sch}
% % \icmlauthor{Firstname8 Lastname8}{yyy,comp}
% % %\icmlauthor{}{sch}
% % %\icmlauthor{}{sch}
\end{icmlauthorlist}

% % \icmlaffiliation{yyy}{Department of XXX, University of YYY, Location, Country}
% % \icmlaffiliation{comp}{Company Name, Location, Country}
% % \icmlaffiliation{sch}{School of ZZZ, Institute of WWW, Location, Country}

\icmlcorrespondingauthor{Firstname1 Lastname1}{first1.last1@xxx.edu}
% % \icmlcorrespondingauthor{Firstname2 Lastname2}{first2.last2@www.uk}

% % % You may provide any keywords that you
% % % find helpful for describing your paper; these are used to populate
% % % the "keywords" metadata in the PDF but will not be shown in the document
% % \icmlkeywords{Machine Learning, ICML}

% % \vskip 0.3in
]

% this must go after the closing bracket ] following \twocolumn[ ...

% This command actually creates the footnote in the first column
% listing the affiliations and the copyright notice.
% The command takes one argument, which is text to display at the start of the footnote.
% The \icmlEqualContribution command is standard text for equal contribution.
% Remove it (just {}) if you do not need this facility.

%\printAffiliationsAndNotice{}  % leave blank if no need to mention equal contribution
% \printAffiliationsAndNotice{\icmlEqualContribution} % otherwise use the standard text.

\begin{abstract}

Parameter-efficient continual learning (PECL) adapts large pre-trained models to evolving task sequences by tuning lightweight adapters while freezing the backbone.
In this setup, cross-task interference is concentrated within the shared adapter subspace. Conventional flatness-based approaches, such as Flat-LoRA, advocate for full-parameter perturbation to achieve global flatness—a requirement at odds with the fixed-backbone design of PECL.
% In this setting, all cross-task interference is funneled into a shared trainable adapter subspace, traditional viewpoints (such as Flat-LoRA) suggest that full-parameter pertur-
% bation is needed to pursue global flatness. This view conflicts with the core design of PECL, where backbone weights are intentionally kept fixed.
To resolve this, we investigate whether flatness restricted to the adapter subspace alone can effectively balance stability and plasticity in PECL. Theoretically, we derive a PAC-Bayesian generalization bound for PECL under a frozen backbone and Gaussian posteriors over the adapter parameters. In this formulation, both the sharpness term and the KL divergence depend solely on the adapter subspace, establishing adapter flatness as the pertinent measure of flatness in PECL.
Experiments on continual learning benchmarks demonstrate that sharpness-aware optimization confined to the adapter subspace consistently enhances performance, whereas full-parameter perturbations yield no systematic improvement. These findings provide a principled foundation for pursuing subspace-restricted flatness in PECL and inform geometry-aware designs for LoRA-based continual learning methods.

% We address this tension by asking whether flatness confined to the adapter subspace is sufficient to govern stability and plasticity. First, we derive a PAC-Bayesian generalization bound specialized to PECL with a frozen backbone and Gaussian posteriors supported on the adapter subspace. In this regime, both the sharpness term and the KL complexity depend only on the adapter subspace, which identifies adapter flatness as the relevant notion of flatness for PECL. 
% % We then relate factor-space flatness to the spectrum of a restricted Fisher information matrix. Finally, 
% Experiments on CL benchmarks show that subspace-restricted sharpness-aware optimization consistently improves continual learning performance, while full-parameter perturbations fail to provide systematic benefits. This provides a principled justification for flatness oriented optimisation confined to the adapter subspace and offers a conceptual basis for geometry aware designs of LoRA based continual learning methods. 


\end{abstract}



\section{Introduction}

Continual learning (CL) has become increasingly relevant for adapting large-scale pre-trained models to evolving tasks in dynamic environments~\cite{PARISI201954,10444954,a2923225ea7d94e6c86faff97e23c69be,10.1145/3735633}. While full fine-tuning remains the most expressive adaptation strategy, it is computationally demanding and highly susceptible to catastrophic forgetting, since newly acquired knowledge may interfere with previously learned representations~\cite{howard2018universallanguagemodelfinetuning,doi:10.1073/pnas.1611835114,pmlr-v234-zhai24a,luo2025empiricalstudycatastrophicforgetting}. Parameter-efficient continual learning (PECL)~\cite{pmlr-v97-houlsby19a,dingParameterefficientFinetuningLargescale2023,He_2025_CVPR,NEURIPS2024_b0bc711f} mitigates these issues by freezing a large pre-trained backbone and confining adaptation to lightweight modules. Low-rank adaptation (LoRA)~\cite{hu2022lora} is a prominent instance of this paradigm, where most task-specific updates are funneled through shared low-rank adapters while the backbone remains frozen.



In PECL, this design has an important consequence. Because the backbone and previously learned adapters are kept fixed, all cross-task interference is concentrated in a shared trainable subspace. Beyond parameter efficiency, a central challenge is therefore to understand and control the local geometry of this adapter subspace, since the shape of the loss landscape in this region governs how new gradients interact with existing adapters and thus mediates the stability–plasticity trade-off.

Flat minima in the loss landscape are widely believed to promote generalization and robustness, since flatter solutions are less sensitive to parameter perturbations~\citep{FlatMinima1997,wu2017understandinggeneralizationdeeplearning,chaudhariEntropySGDBiasingGradient2019,NEURIPS2020_518a38cc,NEURIPS2021_995f5e03,JMLRAnEmpiricalInvestigatio,yangRevisitingFlatnessAwareOptimization2025}.
Sharpness-aware minimization (SAM)~\cite{foretSharpnessAwareMinimizationEfficiently2021} and related methods~\citep{pmlr-v151-bisla22a,zhangGradientNormAware2023,NEURIPS2024_C-Flat} operationalize this by optimizing for parameters that exhibit low loss under small perturbations. 
% Sharpness-aware optimizers such as SAM~\citep{foretSharpnessAwareMinimizationEfficiently2021}, RWP~\citep{pmlr-v151-bisla22a}, GAM~\citep{zhangGradientNormAware2023} and C-FLAT~\cite{NEURIPS2024_C-Flat} instantiate this idea by explicitly penalizing increases of the loss under small parameter perturbations. 
While these methods have proven effective in full fine-tuning, their behavior in the restricted optimization landscape of PECL remains under-explored. 
% LoRA-SAM~\citep{NEURIPS2024_4eb2c0ad} provides evidence that limiting sharpness-aware updates to the LoRA coordinates already yields substantial implicit regularization, whereas Flat-LoRA~\citep{li2025flatloralowrankadaptationflat} argues that flatness restricted to the LoRA subspace may be insufficient and therefore advocates perturbing all model parameters. In the PECL regime that we consider, however, the core efficiency principle is to strictly freeze the pre-trained backbone and confine optimization to the adapter subspace $\mathcal W_\Delta$. This tension raises a central question:
% While effective in full fine-tuning, their application in PECL is not straightforward. 
Recent work presents conflicting views: LoRA-SAM~\citep{NEURIPS2024_4eb2c0ad} suggests that optimizing flatness within the LoRA subspace alone provides effective regularization, while Flat-LoRA~\citep{li2025flatloralowrankadaptationflat} argues that achieving sufficient flatness requires perturbing all model parameters, which implicitly includes the frozen backbone. The latter requirement, however, contradicts the fundamental PECL constraint of a fixed backbone, thus motivating our central question:


% \emph{In PECL, is it necessary to pursue global flatness by perturbing all model parameters, or is controlling flatness solely within the shared adapter subspace sufficient to govern stability and plasticity?}

\emph{In PECL with a frozen backbone, is controlling flatness within the adapter subspace sufficient to govern stability and plasticity, or is global flatness via full-parameter perturbation necessary?}




An initial empirical observation informs our investigation. Applying SAM solely to the adapter parameters can yield solutions residing in broader, flatter basins of the full-model loss landscape (Fig.~\ref{fig:3d_loss_landscape}). This suggests that regularizing the geometry of the adapter subspace can implicitly shape the effective geometry experienced by the model. We therefore hypothesize that subspace-restricted flatness may be adequate to improve robustness and balance in PECL.


% Our investigation is grounded in a simple empirical observation. When SAM is applied only to the low-rank adapter subspace, the loss landscape of the full model around the final solution already becomes wider and smoother than under SGD, even when probed in directions that extend beyond $\mathcal W_\Delta$ (Figure~\ref{fig:3d_loss_landscape}). This suggests that the geometry of the frozen backbone and that of the trainable adapters are not independent. Stabilizing the local curvature within the trainable subspace can implicitly regularize the effective geometry of the whole model. Motivated by this phenomenon, we hypothesize that subspace flatness, understood as geometric control confined to the trainable adapter subspace, may already suffice to substantially improve robustness and achieve a better balance between stability and plasticity in PECL.

% To investigate this hypothesis, we develop a joint  theoretical and empirical analysis. Theoretically, we derive a PAC-Bayesian generalization bound specialized to PECL with a frozen backbone and Gaussian posteriors supported on the adapter subspace.
% In this bound, both the sharpness and complexity terms depend exclusively on the adapter subspace $\mathcal{W}_\Delta$, formally identifying adapter flatness as the relevant notion of flatness for PECL. 
% Our analysis departs from classical hierarchical PAC-Bayesian bounds in two structural ways that are crucial for LoRA-based PECL. First, we derive a task-level CL bound with a time-varying hyper-posterior adapted to the task filtration, which removes the independence assumption on task priors and yields a complexity term that separates task-wise KL divergences from a hyper-posterior drift term capturing cross-task interference. Second, we specialise this bound to PECL with a frozen backbone and a shared LoRA subspace by restricting the posterior support to $\mathcal W_\Delta$, so that both the KL and the empirical terms can be rewritten in terms of adapter-space quantities and an expected sharpness term defined only along $\mathcal W_\Delta$. This subspace-aware formulation is absent from prior PAC-Bayesian analyses and is what allows us to formally distinguish global versus adapter-subspace flatness in the PECL setting.


To investigate this hypothesis, we develop a joint theoretical and empirical analysis. On the theory side, we derive a PAC-Bayesian generalization bound tailored to parameter-efficient continual learning with a frozen backbone and Gaussian posteriors confined to the shared adapter subspace. In contrast to previous PAC-Bayesian bounds, our formulation (i) uses a task-level, time-varying hyper-posterior that relaxes the independence assumption on task priors and separates per-task divergence from a hyper-posterior drift term that measures cross-task interference, and (ii) specializes this bound to the LoRA setting so that both the complexity term and the sharpness-dependent term are expressed purely in adapter-space quantities. This subspace-aware view is absent from prior PAC-Bayesian analyses and is what allows us to formally distinguish global versus adapter-subspace flatness in the PECL regime.


We further connect flatness in the factor space of LoRA to the eigenspectrum of a restricted Fisher information matrix, linking optimization geometry to feature-space stability. Empirically, we compare subspace-restricted and full-parameter sharpness-aware optimization across vision CL benchmarks. 

% Our results show that flatness-aware optimization confined to the adapter subspace consistently improves performance, whereas perturbing all parameters fails to yield systematic gains.

% Under this modeling assumption, the complexity and flatness contributions to the bound can be expressed entirely in terms of quantities restricted to update subspace.



% framework that connects probabilistic generalization guarantees with geometric stability in LoRA-based PECL. On the theoretical side, we derive a PAC-Bayesian generalization bound specialized to PECL with a frozen backbone and Gaussian posteriors supported on the adapter subspace. Under this modelling assumption, both the sharpness term and the KL complexity that govern the bound can be written entirely in terms of quantities restricted to $\mathcal W_\Delta$ rather than to the full parameter space. We further clarify the mechanism linking weight-space flatness to feature-space robustness by showing that LoRA updates implement an abstract update subspace and that optimizing flatness in factor coordinates corresponds to suppressing dominant eigenvalues of the local Fisher information matrix, which stabilizes the learned representations. On the empirical side, we conduct a systematic study on vision benchmarks that compares flatness-aware optimization restricted to $\mathcal W_\Delta$ with full-parameter perturbations within the same PECL models.



Our main contributions are
\begin{itemize}
\item \textbf{PAC-Bayesian Analysis for PECL.} We provide a generalization bound for PECL where the sharpness and complexity are defined over the adapter subspace, offering a theoretical justification for subspace-restricted flatness optimization.
\item \textbf{Geometry of Adapter Flatness.} We show that LoRA updates implement an abstract update subspace and that minimizing curvature in factor space is equivalent to compressing the Fisher spectrum along this subspace. This spectral effect stabilizes class prototypes and feature representations.
\item \textbf{Empirical Evaluation.} On ImageNet R and Tiny ImageNet C/P with ViT B/16 backbones and three adapter organizations (SeqLoRA, IncLoRA, OLoRA), we find that sharpness aware optimization restricted to the adapter subspace consistently improves CL metrics and robustness, whereas full parameter perturbations do not yield systematic gains in the PECL regime and can introduce additional instability.
\end{itemize}

\begin{figure}[tbp]
\centering
\begin{minipage}{1\linewidth}
\centering
\includegraphics[width=1\linewidth]{figures_TRML/2D_lossLandscape_task0_lora_opt_sam_vs_sgd_3d_compare_withxyz.pdf}
\caption{Visualization of the loss landscape for a ViT backbone trained with LoRA adapters, comparing SGD and SAM optimizers. The SAM optimizer not only discovers a wider and flatter basin but also achieves a lower loss value.}
\label{fig:3d_loss_landscape}
\end{minipage}
\end{figure}

% Our main contributions are as follows:
% \begin{itemize}
%     \item \textbf{Sharpness-aware PAC-Bayes for PECL.} We derive a PAC-Bayesian generalization bound tailored to PECL with a frozen backbone and adapter-subspace posteriors. In this regime, both the sharpness and the KL complexity that govern the bound are determined by the trainable adapter subspace $\mathcal W_\Delta$, which identifies adapter-subspace flatness as the relevant notion of flatness for PECL.
%     \item \textbf{Geometric link between adapter flatness and feature stability.} We clarify how LoRA updates realize an abstract update subspace and show that minimizing curvature in factor space is equivalent to suppressing dominant eigenvalues of the restricted Fisher information matrix. This spectral effect stabilizes feature representations and explains how subspace flatness shapes the global loss landscape.
%     \item \textbf{Empirical analysis of flatness strategies in PECL.} We provide systematic evidence on ImageNet-R, ImageNet-C and Tiny-ImageNet-P with ViT-B/16 adapters under three adapter-subspace sharing strategies (SeqLoRA, IncLoRA and OLoRA). Flatness-aware optimization confined to $\mathcal W_\Delta$ consistently improves continual learning metrics and feature stability compared with SGD, whereas full-parameter perturbations do not yield consistent gains in this regime and can introduce additional instability.
% \end{itemize}


\section{Related Work}

\subsection{LoRA-Based PECL and Adapter Subspace Sharing}

Recent work on LoRA-based PECL trains low-rank adapters across tasks over a frozen backbone. A central design axis is the organization of these low-rank adapters within the parameter space, which governs subspace sharing, potential for knowledge transfer, and inter-task interference.
From this viewpoint, LoRA-based PECL methods fall into two regimes: isolated per-task adaptation without subspace interaction, and shared-subspace adaptation where new tasks are trained on top of previously LoRA parameters.


Some methods train an independent LoRA adapter for each task and only activate the corresponding adapter at inference, so there is no parameter sharing and no interaction in the LoRA subspace~\citep{Yu_2024_CVPR}. This almost removes cross-task interference, but also prevents knowledge transfer and requires explicit task selection at inference time.
Shared-Subspace Adaptation builds a common adapter subspace over the frozen backbone, updated sequentially with incoming tasks. Three representative families illustrate how strongly tasks are coupled inside this subspace.
\begin{itemize}
    \item \textbf{Fully shared subspace (SeqLoRA).}
    Early variants reuse a single low-rank adapter for all tasks. A LoRA pair $(A,B)$ is updated sequentially on a fixed backbone $W_0$. Training and inference stay inside the same low-rank subspace, which maximises sharing but also maximizes tasks interference.
    \item \textbf{Incremental shared subspace (IncLoRA and gated variants).}
    Incremental designs expand the adapter subspace over time. IncLoRA introduces a new pair $(A_t,B_t)$ at task $t$ while keeping $W_0$ and earlier adapters fixed and composes all adapters at inference.
% , for example
%     $
%         W_t = W_0 + \sum_{i=1}^t B_i A_i.
%     $
    % The effective update still lies in a common span, so tasks interact inside a shared low-rank subspace, but the per-task decomposition provides partial separation. 
    Gated variants further select or combine adapter with different weight, aiming to reuse useful directions while avoiding unnecessary interference~\citep{Gao_2023_ICCV,10.1145/3730436.3730456,zhao-etal-2024-sapt,cheng2025exploiting,wu2025sdlorascalabledecoupledlowrank}.
    \item \textbf{Constraint-augmented incremental sharing (OLoRA and gradient-level designs).}
    Constraint-augmented methods regularise the geometry of the incremental subspace to mitigate conflicts. OLoRA~\citep{wang2023orthogonal} adds orthogonality penalties to encourage task-specific directions and reduce paramters overlap. Gradient-level approaches~\citep{Liang_2024_CVPR,yang2024parametercollisionhinderingcontinual} align or project gradients away from previously occupied directions. These methods still share a global LoRA subspace, but seek to organise it so that shared adapter capture reusable knowledge while conflicting directions are suppressed.
\end{itemize}




\subsection{Flat Minima}
Prior work has explored flat minima in CL through sharpness-aware optimization~\cite{pmlr-v151-bisla22a,foretSharpnessAwareMinimizationEfficiently2021,zhangGradientNormAware2023}, mode connectivity~\cite{izmailov2019averagingweightsleadswider,mirzadeh2020linearmodeconnectivitymultitask}, and improved representation learning~\cite{hu2022doesselfsupervisedpretrainingperform,ramasesh2022effect}.
These methods are mainly developed for full model training on simple tasks and do not account for the adapter subspace constraints\cite{yangRevisitingFlatnessAwareOptimization2025}. Whether sharpness within the adapter subspace alone can ensure generalization in LoRA-based CL remains unclear, motivating further investigation.


% % Flat minima are widely believed to promote generalization and robustness by reducing the sensitivity of the loss to parameter perturbations~\citep{FlatMinima1997,wu2017understandinggeneralizationdeeplearning,chaudhariEntropySGDBiasingGradient2019,NEURIPS2020_518a38cc,bahri2022sharpnessawareminimizationimproveslanguage,NEURIPS2021_995f5e03,JMLRAnEmpiricalInvestigatio,yangRevisitingFlatnessAwareOptimization2025}. Sharpness-aware optimization methods such as SAM and its variants explicitly minimize adversarial or stochastic sharpness and have been shown to improve both accuracy and robustness in full finetuning regimes~\citep{foretSharpnessAwareMinimizationEfficiently2021,pmlr-v151-bisla22a,Zhang_GAM}. 

% In continual learning, flatness has been studied through mode connectivity and representation smoothing~\citep{izmailov2019averagingweightsleadswider,mirzadeh2020linearmodeconnectivitymultitask,ramasesh2022effect,hu2022doesselfsupervisedpretrainingperform}, but these works operate on the full parameter space and do not distinguish between backbone and adapter directions. 

% More recently, Flat-LoRA-type approaches apply perturbations to the entire parameter space when training LoRA adapters, arguing that subspace flatness may be misaligned with full-space flatness. However, in PECL with a strictly frozen backbone, full-model perturbations conflict with the design principle that old knowledge is preserved in $W_0$. Existing analyses therefore leave open a central question that motivates our work: in LoRA-based PECL, is it sufficient to control flatness within the shared adapter subspace, or must one still reason in the full parameter space to explain forgetting and generalization?

\subsection{PAC-Bayesian Theory }


The PAC-Bayesian framework provides a principled foundation for analyzing generalization by balancing empirical risk and posterior complexity~\cite{mcallesterPACBayesianStochasticModel2003,pentinaPACBayesianBoundLifelong2014,germainPACBayesianTheoryMeets2016,lotfiPACBayesCompressionBounds2022}.
Building on classical PAC-Bayesian theory, hierarchical bounds for lifelong and meta-learning~\citep{pentinaPACBayesianBoundLifelong2014,JMLR:v24:22-1254} learn a hyper-posterior over task-specific priors, but they usually assume that all priors are independent draws from a fixed hyper-prior and therefore do not model the sequential dependence and information accumulation that characterise continual learning. In parallel, online PAC-Bayesian work~\citep{NEURIPS2022_onlineBayes,NEURIPS2024_xinrui}introduces time-varying, data-dependent priors with guarantees that depend on cumulative data, yet it is formulated at the level of individual examples rather than tasks and does not consider task-conditioned posteriors.

% For meta-learning, hierarchical PAC-Bayes bounds introduce a hyper-prior over task-specific priors and derive complexity terms that couple task-level posteriors through the hyper posterior~\citep{pentinaPACBayesianBoundLifelong2014,JMLR:v24:22-1254}. In these formulations, all task priors are assumed to be independent draws from a fixed hyper-prior, which is natural for multi-task transfer but does not encode the sequential dependence and information accumulation that characterize continual learning.
% Parallel work on online PAC-Bayes develops bounds with time-varying, data-dependent priors that are measurable with respect to a filtration of past observations~\citep{NEURIPS2022_onlineBayes,NEURIPS2024_xinrui}. These results capture sequential adaptation at the level of individual examples and have inspired practical algorithms, yet they are not formulated at the task level and remain geometry agnostic: they do not distinguish between different parameter subspaces or relate sharpness to the PAC-Bayesian complexity.

Recent studies connect PAC-Bayes and flat minima by interpreting sharpness as a surrogate for posterior sensitivity~\citep{NIPS2017_10ce03a1,jiang2019fantasticgeneralizationmeasures,mllenhoff2023sam}, which yields bounds for flat solutions in full finetuning. However, these analyses treat all parameters symmetrically and do not address adapter-based PECL, where the posterior is constrained to live in a low-rank subspace on top of a frozen backbone.


% \paragraph{Our position relative to prior PAC-Bayesian work.}


% \paragraph{Our position relative to prior PAC-Bayesian work.}
% Compared to the hierarchical bound of \citet{pentinaPACBayesianBoundLifelong2014}, our analysis makes two structural changes that are essential for LoRA-based PECL. First, we introduce a task-level PAC-Bayesian bound with a time-varying hyper-posterior that is adapted to the task filtration. This removes the independence assumption on task priors and yields a complexity term that decomposes into (i) task-wise KL divergences and (ii) a hyper-posterior drift term that explicitly measures how much the prior over adapters moves across tasks. Ignoring this modification would prevent us from isolating cross-task interference as a separate, quantifiable source of complexity in continual learning.

% Second, we specialize the bound to PECL with a frozen backbone and a shared LoRA subspace. By restricting the posterior support to the adapter subspace $\mathcal W_\Delta$, we show that the full-parameter KL and empirical terms can be rewritten in terms of (i) a KL over adapter updates and (ii) an expected sharpness term that only depends on perturbations within $\mathcal W_\Delta$. This subspace-aware formulation is absent from previous PAC-Bayesian work and is crucial for our main question. Without it, one cannot formally distinguish between global and adapter-subspace flatness, nor explain why LoRA-only sharpness-aware optimization can be beneficial in PECL while full-model perturbations may be systematically harmful. Our experiments later validate these theoretical insights by contrasting local and global perturbations under identical PECL protocols.





\section{Background and Problem Setup}
\label{sec:background}

% We first summarize the continual learning setting, the parameter efficient formulation with shared LoRA adapters, and the notions of flatness and PAC--Bayesian complexity that we use in the rest of the paper.

\subsection{CL and PAC‑Bayesian formulation}
\label{sec:pecl-setup}

Let $\mathcal X$ be the input space, $\mathcal Y$ the label space, and $\mathcal Z = \mathcal X \times \mathcal Y$ the data space.
A model with parameters $W \in \mathbb R^d$ induces a predictor $f_W : \mathcal X \to \mathcal Y$.
We fix a bounded loss $\ell : \mathcal H \times \mathcal Z \to [0,1]$ and write $\ell(W,z)$ for $\ell(f_W,z)$.
In the CL, a learner encounters a sequence of tasks $t\in\{1,\dots,T\}$. 
At task $t$, a dataset
$
S_t = \{z_{tj}\}_{j=1}^{m_t}
$
is drawn i.i.d.\ from a task distribution $\mathcal D_t$.
The empirical risk and population risks on task $t$ are respectively
$
L_{\mathcal D_t}(W)
=
\mathbb E_{z\sim\mathcal D_t}\ell(W,z),
$
$
\hat L_{S_t}(W)
=
\frac{1}{m_t}\sum_{j=1}^{m_t}\ell(W,z_{tj}).
$ 
The PAC-Bayesian viewpoint place probability measures on the hypothesis space rather than focusing on a single function. In this view, for each task $t$, the learner starts from a task-level prior $P_t\in\mathcal M(\mathcal H)$ and, after observing $S_t$, updates it to a task-level posterior $Q_t(\cdot\mid P_t,S_t)\in\mathcal M(\mathcal H)$.
% In this view, learning on a task involves a hypothesis $f \in \mathcal H$ equipped with a prior distribution $P \in \mathcal M(\mathcal H)$, which is updated to a posterior $Q_t(\cdot\mid P_t,S_t)\in\mathcal M(\mathcal H)$ after observing data. 
The corresponding posterior-averaged risks for task $t$ are
% The object of learning on a task has been a hypothesis $f \in \mathcal{H}$ over which one presumes a \textit{prior distribution} $P\sim \mathcal{M}(\mathcal{H})$ that is updated into a \textit{posterior} $Q\sim \mathcal{M}(\mathcal{H})$ based on observed data. We have the corresponding risk as follow
$
L_{\mathcal D_t}(Q_t)
=
\mathbb E_{W_t \sim Q_t(\cdot \mid P_t,S_t)}\mathbb E_{z\sim\mathcal D_t}\ell(W,z),
$ and $
\hat L_{S_t}(Q_t)
=
\mathbb E_{W_t \sim Q_t(\cdot \mid P_t,S_t)}\frac{1}{m_t}\sum_{j=1}^{m_t}\ell(W,z_{tj}).
$


To model the sequential nature of knowledge accumulation across tasks in CL, we adopt a hierarchical PAC‑Bayesian formulation. In this framework, the task-level prior itself is treated as a random variable, and a hyper-posterior is defined as a probability measure $\mathcal P \in \mathcal M(\mathcal M(\mathcal H))$ over such priors~\citep{pentinaPACBayesianBoundLifelong2014,JMLR:v24:22-1254}.
% To analyze the role of sharpness in PECL, we adopt a hierarchical PAC-Bayesian framework, where the prior $P_t$ itself as a random variable.
% Follow hyper PAC , we treat the prior $P_t$ itself as a random variable.
%We write $\mathcal F_t = \sigma(S_1,\dots,S_t)$ for the $\sigma$-algebra generated by the first $t$ datasets and consider the filtration $\mathcal F_0 \subseteq \mathcal F_1 \subseteq \cdots \subseteq \mathcal F_T$. 
% At the beginning of task $t$, before observing $S_t$, the learner maintains a hyper-posterior $\mathcal P_{t-1}$ that is $\mathcal F_{t-1}$-measurable and encodes all information extracted from previous tasks. 
Let $\mathcal F_{t-1}=\sigma(S_1,\dots,S_{t-1})$ denote the $\sigma$-algebra generated by the data observed up to task $t-1$. 
Before observing $S_t$, the learner maintains a hyper-posterior $\mathcal P_{t-1}$ that is $\mathcal F_{t-1}$-measurable, summarizing information transferred from previous tasks. 
% The learning process at task $t$ is then described in two stages. First, a task-level prior $P_t$ is sampled from the current hyper-posterior,
% $
% P_t \sim \mathcal P_{t-1}.
% $
% Second, after observing $S_t$, the learner updates from $P_t$ to a task-level posterior $Q_t$, and the hyper-posterior is updated to $\mathcal P_t$. 
% We assume that these updates are absolutely continuous, so that the associated KL divergences are well defined.
% $
% \mathcal P_t \ll \mathcal P_{t-1},
% \quad
% Q_t(\cdot \mid P,S_t) \ll P(\cdot)
% \quad \text{for } \mathcal P_{t-1}\text{-almost every } P.
% $
% The corresponding task-level risks under the hierarchical posterior are obtained by averaging over the sampling of priors and hypotheses,
The learning procedure at task $t$ proceeds in two stages: (i) sample a task-level prior $P_t$ from the current hyper-posterior, $P_t\sim\mathcal P_{t-1}$; (ii) after observing $S_t$, form the task-level posterior $Q_t$ and update the hyper-posterior to $\mathcal P_t$. 
The corresponding hierarchical risks are obtained by averaging over both the sampled prior and the posterior draw of hypotheses,
$
\mathcal L_t
=
\mathbb E_{P_t \sim \mathcal P_t}
\,
\mathbb E_{W_t \sim Q_t}
\big[
L_{\mathcal D_t}(W_t)
\big],
\hat{\mathcal L}_t
=
\mathbb E_{P_t \sim \mathcal P_t}
\,
\mathbb E_{W_t \sim Q_t}
\big[
\hat L_{S_t}(W_t)
\big].
$


% In PECL, a pre-trained model with frozen weights $W_0$ is adapted to sequential tasks by tuning only lightweight adapter modules. For LoRA, the adapted weights at task $t$ are given by $W_t = W_{\mathrm{frozen}} + \Delta W_t,$ where $W_{\mathrm{frozen}}$ contains the pretrained backbone and all previously learned adapters, which remain frozen at task $t$, and $\Delta W_t$ is a new low rank update confined to a shared adapter subspace $\mathcal W_\Delta$. 
% We focus on PECL methods that reuse a shared adapter subspace on a frozen backbone and distinguish three representative organizations.

% \begin{itemize}
%     \item \textbf{SeqLoRA} (fully shared subspace). A single LoRA pair $(A,B)$ is trained sequentially across tasks, yielding $W_t = W_0 + BA$ at all $t$. Both training and inference remain confined to the same low-rank subspace.
%     \item \textbf{IncLoRA} (incremental partially fused subspaces). Each task $t$ introduces a new pair $(A_t,B_t)$ while keeping $W_0$ and earlier pairs fixed:
%     $
%         W_t = W_0 + \sum_{i=1}^t B_i A_i.
%     $
%     At inference, a fixed or gated combination of adapters is used.
%     \item \textbf{OLoRA} (constraint-augmented sharing). On top of an incremental design, OLoRA~\cite{wang2023orthogonal} introduces explicit orthogonality penalties such as
%     $
%         \sum_{i<t} \|A_i^\top A_t\|_F^2
%     $
%     to encourage task-specific directions within the shared adapter subspace and reduce interference.
% \end{itemize}

% These designs share the same frozen backbone $W_0$ and differ only in how they organize and regularize the shared LoRA subspace $\mathcal W_\Delta$. Methods with completely disjoint adapters that never interact during training or inference are outside our scope.


% We consider three representative strategies for organizing this subspace: a single, fully-shared adapter (SeqLoRA); incrementally added, task-specific adapters (IncLoRA); and adapters regularized for orthogonality to minimize interference (OLoRA).


% In PECL we decompose the parameters at task $t$ as
% $$
% W_t
% =
% W_{\mathrm{frozen}} + \Delta W_t,
% \qquad
% \Delta W_t \in \mathcal W_\Delta \subseteq \mathbb R^d,
% $$
% where $W_{\mathrm{frozen}}$ contains the pretrained backbone and all previously learned adapters, which remain frozen at task $t$, and $\Delta W_t$ is a new low rank update confined to a shared adapter subspace $\mathcal W_\Delta$.

% We focus on three representative organizations of LoRA adapters on the shared backbone.
% SeqLoRA maintains a single adapter in $\mathcal W_\Delta$ and updates it sequentially across tasks.
% IncLoRA accumulates task specific adapters and combines them through a shared head on top of $W_{\mathrm{frozen}}$.
% OLoRA augments this sharing with explicit orthogonality regularization between adapter directions.
% All three methods operate in a common update subspace and differ only in how task specific components are composed inside $\mathcal W_\Delta$.

\subsection{Sharpness Definition}
\label{sec:flatness-subspace}

% Flat minima are often associated with improved generalization and robustness, since they reduce the sensitivity of the empirical loss to small perturbations in parameter space. 
We formalize flatness through sensitivity to parameter perturbations and introduce representative sharpness definitions widely used in the literature.
The adversarial measures the maximal loss increase within a radius $\rho$~\citep{AWP2020,foretSharpnessAwareMinimizationEfficiently2021}: 
$\mathrm{Sh}^{(0)}_S(W,\rho) :=\max_{{\epsilon}\in B(W,\rho)} \Big[\ \hat{L}_{S}(W+{\epsilon})-\hat{L}_{S}(W)\ \Big].$
A %distributional 
stochastic variant averages the loss increase under random perturbations~\citep{liRevisitingRandomWeight2024,jiang2019fantasticgeneralizationmeasures}:
$\mathrm{E\text{-}Sh}_t(\hat{W}_t, \Sigma_t)
 :=
 \mathbb{E}_{\varepsilon \sim \mathcal{N}(0,\Sigma_t)}
 \big[
 \hat{L}_{S_t}(\hat{W}_t+\varepsilon) - \hat{L}_{S_t}(\hat{W}_t)
 \big],$
Typical choices for $\mathcal{N}(0,\Sigma_t)$ include an isotropic Gaussian and a filter-wise Gaussian~\citep{pmlr-v151-bisla22a}. 
% The first order sharpness controls the maximal gradient norm in a neighborhood~\citep{Zhang_GAM}, which is consistent with the zeroth order definition:
% $\mathrm{Sh}^{(1)}_S(W,\rho) := \rho\cdot \max_{W'\in B(W,\rho)} \big\|\nabla_{W} \hat{L}_S(W')\big\|_2,$

Near a local minimum $\hat W$ with $\nabla \hat L_S(\hat W) \approx 0$, a second order expansion gives
$
\hat L_S(\hat W + \varepsilon) - \hat L_S(\hat W)
\approx
\tfrac12\,\varepsilon^\top H(\hat W)\,\varepsilon.
$
The sharpness introduced above reduce to different ways of aggregating the quadratic form $\epsilon^{\top}H(\hat W)\epsilon$: Sh$(0)$ approximates its maximal value over a small norm ball, $\mathrm{E\text{-}Sh}$ approximates its expectation under the perturbation distribution. These identifications justify the use of spectral proxies such as  $\mathrm{tr}\big(H(\hat{W})\big)$, as well as their Fisher surrogates $\mathrm{tr}\big(F(\hat{W})\big)$, where the Fisher coincides with the expected Hessian under standard regularity conditions~\citep{dinhSharpMinimaCan2017,caoTradeoffFlatnessOptimization2024,tsuzukuNormalizedFlatMinima2020,miMakeSharpnessAwareMinimization2022}. 


% If $\varepsilon$ is sampled from $\mathcal N(0,\Sigma)$ inside $\mathcal W_\Delta$, then $\mathrm{E\text{-}Sh}_\Delta(\hat W,\Sigma)$ can be approximated by a quadratic form involving the restriction of $H(\hat W)$ or its Fisher surrogate.

% In this sense, $\mathrm{E\text{-}Sh}_\Delta$ is a curvature dependent notion of flatness that only probes directions along the trainable adapter subspace, which is the relevant geometry in PECL.

% \subsection{A Minimal PAC--Bayes View of Sequential Learning}
% \label{sec:pac-minimal}

% PAC--Bayesian analysis places a prior $P$ and a posterior $Q$ on the hypothesis space and controls generalization through a combination of empirical risk and a KL divergence between $Q$ and $P$.
% Classical PAC-Bayesian results show that, with high probability over the training sample, the expected test risk of a stochastic predictor drawn from $Q$ is upper bounded by its empirical risk plus a complexity term that increases with $\mathrm{KL}(Q\Vert P)$ and decreases with the sample size. Intuitively, good generalization arises when the posterior both fits the data and remains close to the prior.





% In the single task case, for a bounded loss and any prior $P$ and posterior $Q$ with $Q\ll P$, a classical PAC--Bayesian inequality yields, with probability at least $1-\delta$ over the sample $S$,
% $$
% \mathbb E_{W\sim Q} L_{\mathcal D}(W)
% \le
% \mathbb E_{W\sim Q}\hat L_S(W)
% +
% \sqrt{
% \frac{\mathrm{KL}(Q\Vert P) + \log(1/\delta)}
% {2(m-1)}
% }.
% $$
% The right hand side consists of an empirical term and a complexity term that scales with the posterior divergence and the sample size.
% When $Q$ is a Gaussian distribution $Q = \mathcal N(\hat W,\Sigma)$, the empirical term $\mathbb E_{W\sim Q}\hat L_S(W)$ can be written as the loss at $\hat W$ plus an expected sharpness term.

% In Section~\ref{sec:pac-pecl} we build on a task level extension of this inequality with a time varying hyper posterior over task priors.
% This yields a continual learning bound whose complexity decomposes into task wise KL divergences and a hyper posterior drift across tasks.
% We then specialise this bound to the PECL parameterisation above and show that both the empirical term and the complexity term can be expressed entirely in the adapter subspace $\mathcal W_\Delta$.


\section{A Sharpness-Aware PAC-Bayes Bound for PECL}
\label{sec:pac-pecl}


\begin{figure}[thph]
\centering
\begin{minipage}{1\linewidth}
\centering
\includegraphics[width=\linewidth]{figures_TRML/shiyitu668.pdf}
% \caption{Illustration of the KL divergence decomposition and coordinate equivalence. 
%     (1) \textbf{KL Decomposition}: In PECL, the overall model complexity, measured by $\mathrm{KL}(Q^{\mathrm{full}} \| P^{\mathrm{full}})$, is solely determined by the distribution of the updated parameters in the adapter subspace $\mathcal{W}_\Delta$, formalized as $\mathrm{KL}(Q^{\Delta} \| P^{\Delta})$, while the frozen backbone contributes no divergence. 
%     (2) \textbf{Two equivalent views}: The optimization in the low-rank factor coordinates $(A, B)$ is theoretically isomorphic to operating directly in the update subspace $\mathcal{W}_\Delta$. This equivalence preserves geometric properties such as flatness and curvature, which allows sharpness-aware optimization to be confined to the factorized LoRA parameters.}
\vspace{-4pt}
\caption{KL decomposition and coordinate equivalence in PECL. With a frozen backbone, the complexity term satisfies
$\mathrm{KL}(Q^{\mathrm{full}}\|P^{\mathrm{full}})=\mathrm{KL}(Q^{\Delta}\|P^{\Delta})$, so it depends only on the adapter subspace $\mathcal W_{\Delta}$.
LoRA factor coordinates $(A,B)$ parameterize the same update subspace, enabling sharpness control to act directly on the trainable factors.}
\label{fig: KL_decompose}
\end{minipage}
\end{figure}


% In parameter efficient continual learning (PECL) the model at task $t$ is written as
% $$
% W_t \;=\; W_{\mathrm{frozen}} + \Delta W_t,
% $$
% where $W_{\mathrm{frozen}}$ collects the pretrained backbone and all previously learned adapters, and $\Delta W_t$ is a new low rank update confined to a shared adapter subspace $\mathcal W_\Delta \subseteq \mathbb R^d$. This section develops a PAC–Bayesian analysis that isolates the role of \emph{adapter subspace flatness} in PECL.

% We adopt a hierarchical PAC–Bayesian view with task level priors and posteriors sampled from a hyper posterior over tasks. A pathwise PAC–Bayes inequality for continual learning with time varying hyper posteriors is provided in Appendix~A (Theorem~A.1). It bounds the task averaged test risk by an empirical term evaluated under Gaussian perturbations and two KL complexity terms on task posteriors and hyper posteriors. In what follows we specialise this general result to the PECL parameterisation and show that both flatness and complexity can be written entirely in the adapter subspace.

\subsection{KL decomposition under a frozen backbone}
\label{sec:kl-decomp-pecl}

% For each task $t$, the PECL parameterisation writes
% $$
% W_t \;=\; W_{\mathrm{frozen}} + \Delta W_t,
% \qquad
% \Delta W_t \in \mathcal W_\Delta \subseteq \mathbb R^d,
% $$
% where $\mathcal W_\Delta$ is the linear subspace of admissible adapter updates. We can therefore use $\Delta W$ as coordinates on the affine set $W_{\mathrm{frozen}} + \mathcal W_\Delta$.

% Given a prior $P_t^\Delta$ and a posterior $Q_t^\Delta$ on the adapter subspace, we define the corresponding distributions on full parameters simply by translation
% \begin{equation*}
% \begin{aligned}
%         W \sim P_t^{\mathrm{full}}
% \iff
% W = W_{\mathrm{frozen}} + \Delta W,
% \quad \Delta W \sim P_t^\Delta, \\
% \quad
% W \sim Q_t^{\mathrm{full}}
% \iff
% W = W_{\mathrm{frozen}} + \Delta W,
% \quad \Delta W \sim Q_t^\Delta.
% \end{aligned}
% \end{equation*}

Let $\mathcal W_\Delta\subseteq\mathbb R^d$ be a linear subspace and let $W+\mathcal W_\Delta$ be its affine translate. Equip both sets with the trace Borel $\sigma$-algebra induced from $\mathbb R^d$. $
\mathcal B(\mathcal W_\Delta)\ :=\{B\cap \mathcal W_\Delta:\ B\in \mathcal B(\mathbb R^d)\},\quad
\mathcal B(W+\mathcal W_\Delta)\ :=\{B\cap (W+\mathcal W_\Delta):\ B\in \mathcal B(\mathbb R^d)\}.
$
These coincide with the Borel $\sigma$-algebras generated by the subspace topologies on $\mathcal W_\Delta$ and $W+\mathcal W_\Delta$, respectively, so both are standard Borel spaces.


To formalize the probabilistic relationship between the LoRA subspace and the full model space, we define the translation map:
% \begin{align*}
$T_W: (\mathcal W_\Delta,\mathcal B(\mathcal W_\Delta)) \to (W+\mathcal W_\Delta,\mathcal B(W+\mathcal W_\Delta)), 
T_W(\Delta)=W+\Delta,$
% \end{align*}
which is a Borel isomorphism between the subspace and affine space.  The full-space prior and posterior are defined as pushforwards onto the affine subspace by
$
P^{\mathrm{full}}:=(T_W)_{\text{\#}} P^{\Delta},
Q^{\mathrm{full}}:=(T_W)_{\text{\#}} Q^{\Delta},
$
where $P_t^{\Delta}$ and $Q_t^{\Delta}$ are probability measures on $(\mathcal{W}\Delta, \mathcal{B}(\mathcal{W}\Delta))$.
% This construction is introduced to ensure that the PAC-Bayesian complexity term in the full parameter space can be expressed entirely in the adapter subspace under the PECL constraint. 
% For any measurable $F$, the pushforward $F_{\text{\#}}\mu$ is defined by $(F_{\#}\mu)(A)=\mu(F^{-1}(A))$. 
% This pushforward construction leads to a fundamental simplification of the KL divergence between full-model distributions.


% , as established in the following lemma.
We propose the following lemma to formalize the resulting translation-invariance of the KL divergence, and its proof is included in Appendix~\ref{proof:Lemma1}.
\begin{lemma}[KL in PECL is adapter-subspace KL]
\label{lem:kl-decomp-pecl}
Assume $Q_t^\Delta \ll P_t^\Delta$. Then $Q_t^{\mathrm{full}} \ll P_t^{\mathrm{full}}$ and
$$
\mathrm{KL}\bigl(Q_t^{\mathrm{full}} \,\Vert\, P_t^{\mathrm{full}}\bigr)
\;=\;
\mathrm{KL}\bigl(Q_t^\Delta \,\Vert\, P_t^\Delta\bigr).
$$
\end{lemma}
The full model distributions $P_t^{\mathrm{full}}$ and $Q_t^{\mathrm{full}}$ are pushforwards of $P_t^\Delta$ and $Q_t^\Delta$ under the affine map $\Delta W \mapsto W_{\mathrm{frozen}} + \Delta W$. Since this map is bijective on $W_{\mathrm{frozen}} + \mathcal W_\Delta$ and preserves Radon–Nikodym derivatives, the relative entropy is invariant. In PECL, the KL complexity is therefore completely determined by the distribution of the adapter update $\Delta W\in\mathcal W_\Delta$; the frozen backbone contributes no divergence.

\subsection{Expected sharpness restricted to the adapter subspace}
\label{sec:subspace-sharpness-pecl}

We specialise the posterior to a Gaussian supported on $\mathcal W_\Delta$. For task $t$ We consider a Gaussian posterior \( Q_t^{\Delta} = \mathcal{N}(\hat{\Delta W}_t, \Sigma_t) \) on the measurable space \( (\mathcal{W}_\Delta, \mathcal{B}(\mathcal{W}_\Delta)) \), where \( \hat{\Delta W}_t \in \mathcal{W}_\Delta \) is the mean and \( \Sigma_t\) is a covariance on the adapter subspace \( \mathcal{W}_\Delta \). Sampling from \( \Delta W \sim Q_t^{\Delta} \) is equivalent to 
$
\Delta W = \hat{\Delta W}_t + \varepsilon, \quad \varepsilon \sim \mathcal{N}(0, \Sigma_t) \ \text{on } \mathcal{W}_\Delta.
$ This Gaussian posterior formulation enables a decomposition of the empirical risk into nominal performance and subspace-restricted sharpness.


We propose the following lemma for this decomposition, and its proof is included in Appendix~\ref{proof:lemma2}.
\begin{lemma}[Empirical loss and subspace sharpness]
\label{lem:subspace-loss-pecl}
For each task $t$,
\begin{equation*}
\label{eq:subspace-loss-pecl}
\begin{aligned}
&\mathbb E_{W\sim Q_t^{\mathrm{full}}}\bigl[\hat L_{S_t}(W)\bigr] \\
&= \mathbb{E}_{\varepsilon\sim\mathcal N(0,\Sigma_t)}
    \big[\hat{L}_{S_t}(W_\mathrm {frozen }+\hat{\Delta W}_t+\varepsilon)\big],\varepsilon\in\mathcal W_\Delta \\
& =\hat L_{S_t}(W_{\mathrm{frozen}} + \hat{\Delta W}_t)
+
\mathrm{E\text{-}Sh}_t^\Delta(\hat{\Delta W}_t,\Sigma_t),\varepsilon\in\mathcal W_\Delta.
\end{aligned}
\end{equation*}
where the expected sharpness $
\mathrm{E\text{-}Sh}_t^{\Delta}(\hat{\Delta W}_t,\Sigma_t)
:= \mathbb{E}_{\varepsilon\sim\mathcal N(0,\Sigma_t)}
   \big[\hat{L}_{S_t}(W_\mathrm {frozen }+\hat{\Delta W}_t+\varepsilon)
       - \hat{L}_{S_t}(W_\mathrm {frozen }+\hat{\Delta W}_t)\big],\quad
\varepsilon\sim\mathcal N(0,\Sigma_t) \in  \mathcal W_\Delta$
% \begin{equation*}
% \begin{aligned}
% &\mathrm{E\text{-}Sh}_t^{\Delta}(\hat{\Delta W}_t,\Sigma_t):= \\
% &\qquad \mathbb{E}_{\varepsilon\sim\mathcal N(0,\Sigma_t)}
%    \big[\hat{L}_{S_t}(W_\mathrm {frozen }+\hat{\Delta W}_t+\varepsilon) \\
% & \qquad- \hat{L}_{S_t}(W_\mathrm {frozen }+\hat{\Delta W}_t)\big] \\
% &\qquad \varepsilon\sim\mathcal N(0,\Sigma_t) \in \mathcal W_\Delta
% \end{aligned}
% \end{equation*}
\end{lemma}
Lemma~\eqref{lem:subspace-loss-pecl} decomposes the posterior expected empirical loss into the empirical loss at the mean adapter plus an increase that depends only on perturbations inside $\mathcal W_\Delta$. The backbone is frozen and does not appear in the sharpness term. In PECL, the relevant notion of flatness is therefore fully captured by sharpness restricted to the adapter subspace.





\subsection{Sharpness-aware PAC–Bayes bound for PECL}
\label{sec:pecl-bound-main}




% We combine the hierarchical PAC-Bayes with Lemma~\ref{lem:kl-decomp-pecl} and Lemma~\ref{lem:subspace-loss-pecl} . This yields a sharpness-aware bound whose empirical and complexity terms are both confined to the adapter subspace. The proof is in Appendix~\ref{proof:theorem}.


% We combine the hierarchical PAC-Bayes framework with Lemma~\ref{lem:kl-decomp-pecl} and Lemma~\ref{lem:subspace-loss-pecl} to obtain the following sharpness-aware PAC-Bayesian bound, in which both the empirical term and the complexity term are confined to the adapter subspace. We below provide the theorem\ref{thm:pecl-bound}, and its proof is included in Appendix~\ref{proof:theorem}.

We combine the hierarchical PAC Bayes framework with Lemma~\ref{lem:kl-decomp-pecl} and Lemma~\ref{lem:subspace-loss-pecl} to derive a sharpness aware bound whose empirical and complexity terms are both confined to the adapter subspace.
Lemma~\ref{lem:kl-decomp-pecl} is illustrated in Fig.~\ref{fig: KL_decompose}, showing that the full space divergence reduces to $\mathrm{KL}(Q^{\Delta}\|P^{\Delta})$ under a frozen backbone, while Lemma~\ref{lem:subspace-loss-pecl} expresses the empirical term through adapter space perturbations.
Fig.~\ref{fig: showPAC3} summarizes the resulting hierarchical construction across tasks and the corresponding decomposition of the bound.
We now state Theorem~\ref{thm:pecl-bound}; the proof is provided in Appendix~\ref{proof:theorem}.


% [Sharpness-aware PAC--Bayes bound for PECL]
\begin{theorem}
\label{thm:pecl-bound}
Consider a PECL model with frozen backbone $W_{\mathrm{frozen}}$ and Gaussian posteriors 
$Q_t^\Delta = \mathcal N(\hat{\Delta W}_t,\Sigma_t)$ supported on $\mathcal W_\Delta$. 
Let $P_t^\Delta$ be task priors sampled from a hyper posterior $\mathcal P_{t-1}$ over adapter 
distributions. For any $\delta\in(0,1]$, with probability at least $1-\delta$ over $S_{1:T}$,
\begin{equation}
\label{eq:pecl-main-bound}
\resizebox{\linewidth}{!}{$
\begin{aligned}
&\frac{1}{T}\sum_{t=1}^T
\mathbb E_{P_t^\Delta\sim\mathcal P_{t-1}}
\mathbb E_{W\sim Q_t^{\mathrm{full}}}
L_{\mathcal D_t}(W) \\
&\le
\frac{1}{T}\sum_{t=1}^T
\mathbb E_{P_t^\Delta\sim\mathcal P_{t-1}}
\bigl[
\hat L_{S_t}(W_{\mathrm{frozen}} + \hat{\Delta W}_t)
+
\mathrm{E\text{-}Sh}_t^\Delta(\hat{\Delta W}_t,\Sigma_t)
\bigr] \\
&\quad+
\sqrt{
\frac{
\displaystyle
\sum_{t=1}^T
\mathbb E_{P_t^\Delta\sim\mathcal P_{t-1}}
\mathrm{KL}\bigl(Q_t^\Delta \,\Vert\, P_t^\Delta\bigr)
+
\sum_{t=1}^T
\mathrm{KL}\bigl(\mathcal P_t\Vert\mathcal P_{t-1}\bigr)
+
\log\frac{1}{\delta}}
{2\displaystyle\sum_{t=1}^T (m_t - 1)}
}.
\end{aligned}
$}
\end{equation}
\end{theorem}

Theorem~\ref{thm:pecl-bound} specialises the hierarchical PAC-Bayesian framework to a PECL regime with a frozen backbone and posteriors supported on the adapter subspace $\mathcal{W}_\Delta$. In this setting, the empirical term is evaluated under Gaussian perturbations \( \varepsilon \in \mathcal W_\Delta \), which gives rise to the adapter-subspace expected sharpness \( \mathrm{E\text{-}Sh}^{\Delta}_t(\hat{\Delta W_t},\Sigma_t) \), while the complexity term decomposes into adapter-level divergences \( \mathrm{KL}(Q^\Delta_t \Vert P^\Delta_t) \) and a hyper-posterior drift \( \mathrm{KL}(P_t \Vert P_{t-1}) \). Since the backbone is frozen, neither sharpness nor complexity depends on directions outside \( \mathcal W_\Delta \), so the bound singles out adapter-subspace flatness as the relevant notion of flatness in PECL and theoretically justifies focusing flatness-aware optimisation on the adapter subspace.








% \begin{theorem}[Sharpness-aware PAC–Bayes bound for PECL]
% \label{thm:pecl-bound}
% Consider a PECL model with frozen backbone $W_{\mathrm{frozen}}$ and Gaussian posteriors $Q_t^\Delta = \mathcal N(\hat{\Delta W}_t,\Sigma_t)$ supported on $\mathcal W_\Delta$. Let $P_t^\Delta$ be task priors sampled from a hyper posterior $\mathcal P_{t-1}$ over adapter distributions. For any $\delta\in(0,1]$, with probability at least $1-\delta$ over $S_{1:T}$,
% \begin{equation}
% \label{eq:pecl-main-bound}
% \begin{aligned}
% \frac1T\sum_{t=1}^T
% \mathbb E_{P_t^\Delta\sim\mathcal P_{t-1}}
% \mathbb E_{W\sim Q_t^{\mathrm{full}}}
% L_{\mathcal D_t}(W)
% \;\le\;
% &\;
% \frac1T\sum_{t=1}^T
% \mathbb E_{P_t^\Delta\sim\mathcal P_{t-1}}
% \Bigl[
% \hat L_{S_t}(W_{\mathrm{frozen}} + \hat{\Delta W}_t)
% +
% \mathrm{E\text{-}Sh}_t^\Delta(\hat{\Delta W}_t,\Sigma_t)
% \Bigr]
% \\[0.2em]
% &\;+\;
% \sqrt{
% \frac{
% \displaystyle \sum_{t=1}^T
% \mathbb E_{P_t^\Delta\sim\mathcal P_{t-1}}
% \mathrm{KL}\bigl(Q_t^\Delta \,\Vert\, P_t^\Delta\bigr)
% \;+\;
% \displaystyle \sum_{t=1}^T
% \mathrm{KL}(\mathcal P_t\Vert\mathcal P_{t-1})
% \;+\;
% \log\frac1\delta}
% {2\displaystyle\sum_{t=1}^T (m_t - 1)}
% }.
% \end{aligned}
% \end{equation}
% \end{theorem}

% Theorem~\ref{thm:pecl-bound} decomposes the task averaged test risk into three components.

% \begin{itemize}
% \item An empirical term consisting of the empirical loss at $W_{\mathrm{frozen}}+\hat{\Delta W}_t$ and the adapter-subspace expected sharpness $\mathrm{E\text{-}Sh}_t^\Delta(\hat{\Delta W}_t,\Sigma_t)$.
% \item A sum of adapter-level KL divergences $\sum_t \mathbb E_{P_t^\Delta}\mathrm{KL}(Q_t^\Delta\Vert P_t^\Delta)$, which measures how much each task posterior deviates from its prior inside $\mathcal W_\Delta$.
% \item A hyper-posterior drift $\sum_t \mathrm{KL}(\mathcal P_t\Vert\mathcal P_{t-1})$, which quantifies how the distribution of task priors evolves with the task sequence.
% \end{itemize}

% Curvature enters the bound only through the expected sharpness term, which is defined over Gaussian perturbations supported on $\mathcal W_\Delta$. The complexity term is also expressed entirely through KL divergences on $\mathcal W_\Delta$ and their hyper-level evolution. In a PECL regime, adapter-subspace flatness is therefore the relevant notion of flatness: tightening the bound requires controlling $\mathrm{E\text{-}Sh}_t^\Delta$ and the adapter KL divergences, while directions outside $\mathcal W_\Delta$ play no role.

A special case is SeqLoRA, where the prior is chosen as the previous posterior, $P_t^\Delta := Q_{t-1}^\Delta$ conditioned on the realised history. In this case the hyper update is deterministic, $\mathcal P_t = \mathcal P_{t-1}$ and the hyper drift term vanishes. The complexity reduces to $\sum_t \mathrm{KL}(Q_t^\Delta\Vert Q_{t-1}^\Delta)$, yielding a trade-off between subspace sharpness and the cumulative adapter KL.


\begin{figure}[thph]
\centering
% \begin{minipage}{1\linewidth}
% \centering
\includegraphics[width=\linewidth]{figures_TRML/showPAC3.pdf}
% \caption{Illustration of the PAC bound.}
\label{fig: showPAC3}
% \end{minipage}
\vspace{-10pt}
\caption{Hierarchical PAC Bayes in PECL. A time-varying hyperposterior $\mathcal{P}_t$ induces task-level priors ${P_t}$. The resulting bound consists of an empirical term and a complexity term, where the latter decomposes into the task-wise divergences $\sum_t \mathbb{E}_{\mathcal{P}_t}\mathrm{KL}(Q_t|P_t)$ and the hyperposterior drift $\sum_t \mathrm{KL}(\mathcal{P}_t|\mathcal{P}_{t-1})$.}
\end{figure}

\subsection{Geometric interpretation via LoRA factors and feature Fisher}
\label{sec:geometry-lora-pecl}

The bound in Theorem~\ref{thm:pecl-bound} is formulated on an abstract update subspace $\mathcal W_\Delta$. In practice, this subspace is realised by LoRA style low-rank adapters, and the curvature that controls $\mathrm{E\text{-}Sh}_t^\Delta$ can be expressed in terms of feature-space matrices and the classifier.

Let $\phi(x;W)\in\mathbb R^D$ be the representation produced by the backbone and adapters, and let $W_{\mathrm{cls}}\in\mathbb R^{D\times K}$ be the linear classifier. The predictive distribution is
$
p(y\mid x;W,W_{\mathrm{cls}})
=
\mathrm{softmax}\bigl(W_{\mathrm{cls}}^\top \phi(x;W)\bigr).
$
The local feature matrix at $x$ is
$
E_\phi(x)
:=
\mathbb E_{y\sim p(\cdot\mid x;W,W_{\mathrm{cls}})}
\bigl[
\nabla_\phi \log p
\nabla_\phi \log p^\top
\bigr]. 
$\cite{EFC2024} By the chain rule, the Fisher information matrix for parameters $W$ at $x$ is
$
F_W(x)
=
\Bigl(
\frac{\partial \phi(x;W)}{\partial W}
\Bigr)^\top
E_\phi(x)
\frac{\partial \phi(x;W)}{\partial W}.
$
For a small perturbation $\Delta W$, the change in predictive distribution satisfies the quadratic approximation
$$
\mathrm{KL}\bigl(p(\cdot\mid x;W+\Delta W)\,\Vert\,p(\cdot\mid x;W)\bigr)
\approx
\frac12\Delta W^\top F_W(x)\Delta W.
$$
Restricting to admissible updates $\Delta W\in\mathcal W_\Delta$ yields the curvature that governs the subspace expected sharpness $\mathrm{E\text{-}Sh}_t^\Delta$.

LoRA realises $\mathcal W_\Delta$ by low rank factors. For a layer with weight matrix $W_\ell \in \mathbb R^{d_o^{(\ell)}\times d_i^{(\ell)}}$, a LoRA module introduces an update
$
\Delta W_\ell = B_\ell A_\ell.
$ Let $(\widehat A_\ell,\widehat B_\ell)$ be the current factors and $(\delta A_\ell,\delta B_\ell)$ small perturbations. Collecting all factor perturbations into a vector $\theta$ and linearising gives
$
\mathrm{vec}(\Delta W)
\;\approx\;
J_{\mathrm{LoRA}}\,\theta,
$
where $J_{\mathrm{LoRA}}$ is the Jacobian from factor coordinates to full parameters. The image of $J_{\mathrm{LoRA}}$ coincides with $\mathcal W_\Delta$, so variations of $\theta$ span exactly the update subspace. A Gaussian posterior on the factor coordinates induces a Gaussian posterior on $\Delta W$ supported on $\mathcal W_\Delta$.

This correspondence implies that flatness-aware optimisation applied to LoRA factors, such as SAM or GAM, directly reduces the dominant eigenvalues of the Fisher matrix restricted to $\mathcal W_\Delta$. For a fixed posterior covariance this contraction decreases $\mathrm{E\text{-}Sh}_t^\Delta$ and allows the adapter posterior to remain closer to its prior, which jointly tightens the sharpness term and the KL terms in Theorem~\ref{thm:pecl-bound}. Through the chain rule the same spectral contraction implies that the representation $\phi(x;W)$ becomes less sensitive to LoRA updates.

% Class prototypes drift less, empirical feature matrices exhibit larger effective rank and lower anisotropy, and NME-based evaluation becomes more stable. The empirical trends we observe in AAA, NME, prototype drift, CKA and feature-spectrum metrics are therefore consistent with the mechanism predicted by the bound: adapter-subspace flatness simultaneously improves PAC–Bayesian complexity and feature stability in PECL.

\begin{table*}[tph]
\centering
\caption{
Class incremental performance, weight-space flatness and feature-space robustness on ImageNet-R ($T=20$, $R=16$).
}
\label{tab:table_all}

% 上下空白略微压缩
% \vspace{-0.1em}

% 列间距和行高压缩一点
\setlength{\tabcolsep}{2pt}      % 原来是 3pt
\renewcommand{\arraystretch}{0.85} % 原来是 0.95
\small                            % 字体略小一号

% 如果不想整体缩放，可以去掉 resizebox
\resizebox{0.97\textwidth}{!}{   % 比原来略缩小一点，顺带减小高度
\begin{tabular}{l l cccc ccc cc}
\toprule
\multirow{2}{*}{\textbf{Method}} & \multirow{2}{*}{\textbf{Opt.}} &
\multicolumn{4}{c}{\textbf{CL Performance} ($\uparrow$)} &
\multicolumn{3}{c}{\textbf{Weight Flatness} ($\downarrow$)} &
\multicolumn{2}{c}{\textbf{Feature Robustness}} \\
\cmidrule(lr){3-6}\cmidrule(lr){7-9}\cmidrule(lr){10-11}
& &
$\operatorname{ACC}_{20}^{\mathrm{cls}}$ &
$\operatorname{AAA}^{\mathrm{cls}}$ &
$\operatorname{ACC}_{20}^{\mathrm{ncm}}$ &
$\operatorname{AAA}^{\mathrm{ncm}}$ &
Sh$^{(0)}(\rho)$ 
& $\lambda_{\max}(F)$ &
$tr(F)$ &
L2 drift ($\downarrow$) &
$\mathrm{tr}(E)$ \\
\midrule
\multirow{1}{*}{\textbf{PerTaskLP}}
& SGD
& $55.22^{\pm 0.31}$ & $59.36^{\pm 0.19}$
& $62.96^{\pm 0.12}$ & $66.58^{\pm 0.37}$
& 0.072  & 8.47  & 1015.07  
& 0.00 & 0.82 \\
\midrule
\multirow{3}{*}{\textbf{SeqFT}}
& SGD
& $50.72^{\pm 2.20}$ & $59.89^{\pm 1.32}$
& $69.49^{\pm 0.16}$ & $75.10^{\pm 0.57}$
& 0.582   & 538.79               & 29703.66
& 29.57 & 0.58 \\
& SAM
& $56.07^{\pm 2.60}$ & $66.47^{\pm 1.53}$
& $72.36^{\pm 2.46}$ & $78.31^{\pm 0.50}$
& 0.259 & 75.43                & 5881.32
& 17.04 & 0.74 \\
& GAM
& $\mathbf{59.78}^{\pm 0.74}$ & \textbf{$\mathbf{67.31}^{\pm 1.18}$}
& $\mathbf{73.74}^{\pm 0.88}$ & \textbf{$\mathbf{78.32}^{\pm 0.94}$}
& 0.177 & 32.00                & 3273.08      
& 13.80 & 0.78 \\
\midrule
\multirow{3}{*}{\textbf{SeqLoRA}}
& SGD
& $63.14^{\pm 2.17}$ & $69.07^{\pm 1.33}$
& $72.82^{\pm 1.27}$ & $76.62^{\pm 0.68}$
& 0.142 & 35.51                & 1401.15
& 16.58 & 0.77 \\
& SAM
& $67.04^{\pm 0.71}$ & $71.42^{\pm 0.71}$
& $74.53^{\pm 0.60}$ & $77.62^{\pm 0.43}$
& 0.120 & 21.39                & 821.45
& 14.82 & 0.81 \\
& GAM
& $\mathbf{73.31}^{\pm 0.52}$ & \textbf{$\mathbf{76.80}^{\pm 0.56}$}
& $\mathbf{75.74}^{\pm 0.15}$ & \textbf{$\mathbf{78.86}^{\pm 0.53}$}
& 0.060 & 3.72                 & 240.91 
& 11.28 & 0.78 \\
\midrule
\multirow{3}{*}{\textbf{IncLoRA}}
& SGD
& $65.77^{\pm 2.84}$ & $71.29^{\pm 0.78}$
& $75.41^{\pm 0.26}$ & $78.37^{\pm 0.26}$
& 0.076 & 8.89                 & 836.91    
& 16.87 & 0.67 \\
& SAM
& $68.86^{\pm 0.92}$ & $72.59^{\pm 0.39}$
& $75.95^{\pm 0.07}$ & $78.68^{\pm 0.25}$
& 0.053 & 4.35                 & 526.49     
& 15.14 & 0.70 \\
& GAM
% & $72.86^{\pm 0.23}$ & \textbf{$76.52^{\pm 0.48}$}
% & $76.72^{\pm 0.54}$ & \textbf{$79.42^{\pm 0.40}$}
& $\mathbf{72.86}^{\pm 0.23}$ & $\mathbf{76.52}^{\pm 0.48}$
& $\mathbf{76.72}^{\pm 0.54}$ & $\mathbf{79.42}^{\pm 0.40}$ 
& 0.035 & 2.09                 & 192.25 
& 12.15 & 0.73 \\
\midrule
\multirow{3}{*}{\textbf{OLoRA}}
& SGD
& $66.64^{\pm 0.44}$ & $71.41^{\pm 0.47}$
& $75.10^{\pm 0.41}$ & $78.65^{\pm 0.67}$
& 0.078 & 9.00                 & 831.23
& 17.23 & 0.68 \\
& SAM
& $68.53^{\pm 0.78}$ & $72.75^{\pm 0.76}$
& $75.73^{\pm 0.32}$ & $78.70^{\pm 0.67}$
& 0.066 & 7.41                 & 628.12   
& 15.17 & 0.70 \\
& GAM
% & $72.44^{\pm 0.36}$ & \textbf{$76.05^{\pm 0.52}$}
% & $76.56^{\pm 0.53}$ & \textbf{$79.16^{\pm 0.43}$}
& $\mathbf{72.44}^{\pm 0.36}$ & $\mathbf{76.05}^{\pm 0.52}$
& $\mathbf{76.56}^{\pm 0.53}$ & $\mathbf{79.16}^{\pm 0.43}$
& 0.036& 2.08                 & 198.30  
& 12.16 & 0.74 \\
% \midrule
% \multirow{1}{*}{\textbf{l2p}}
% & Adam
% & $71.35^{\pm 0.05}$ & $74.84^{\pm 0.73}$
% & -- & --
% & -- & --
% & -- & -- \\
% \midrule
% \multirow{1}{*}{\textbf{dualprompt}}
% & Adam
% & $71.59^{\pm 0.68}$ & $74.80^{\pm 1.18}$
% & -- & --
% & -- & --
% & -- & -- \\
\bottomrule
\end{tabular}
}
\end{table*}


\section{Experimental Setup}



% \begin{table*}[htpb]
% \centering
% \caption{
% Class incremental performance, weight-space flatness and feature-space robustness on ImageNet-R ($T=20$, $R=16$). 
% }
% \label{tab:table_all}
% \setlength{\tabcolsep}{3pt}
% \renewcommand{\arraystretch}{0.95}
% \resizebox{\linewidth}{!}{
% \begin{tabular}{l l cccc cc cc}
% \toprule
% \multirow{2}{*}{\textbf{Method}} & \multirow{2}{*}{\textbf{Opt.}} &
% \multicolumn{4}{c}{\textbf{CL Performance} ($\uparrow$)} &
% \multicolumn{2}{c}{\textbf{Weight Flatness} ($\downarrow$)} &
% \multicolumn{2}{c}{\textbf{Feature Robustness}} \\
% \cmidrule(lr){3-6}\cmidrule(lr){7-8}\cmidrule(lr){9-10}
% & &
% $\operatorname{ACC}_{20}^{\mathrm{cls}}$ &
% $\operatorname{AAA}^{\mathrm{cls}}$ &
% $\operatorname{ACC}_{20}^{\mathrm{nme}}$ &
% $\operatorname{AAA}^{\mathrm{nme}}$ &
% Sh$^{(0)}(\rho)$ &
% $\mathrm{tr}(F)$ &
% L2 drift ($\downarrow$) &
% $\mathrm{tr}(E)$ \\
% \midrule



\subsection{Datasets and task protocol}

We evaluate PECL methods on vision benchmarks derived from ImageNet. ImageNet-R~\cite{ImagenetR_Hendrycks_2021_ICCV} serves as our main class-incremental benchmark and is split into $T=20$ tasks over $200$ classes following prior work~\cite{liangInfLoRAInterferenceFreeLowRank2024a}. To assess robustness under corruptions and perturbations, we further use Tiny ImageNet-C and Tiny ImageNet-P~\cite{hendrycks2018benchmarking}. 
% Tiny ImageNet-C applies common corruptions to Tiny ImageNet validation images, while Tiny ImageNet-P provides short perturbation sequences per image.



% \begin{table*}[htpb]
% \centering
% \caption{
% Class incremental performance, weight-space flatness and feature-space robustness on ImageNet-R ($T=20$, $R=16$). 
% }
% \label{tab:table_all}
% \setlength{\tabcolsep}{4pt}
% \renewcommand{\arraystretch}{0.95}
% \resizebox{\linewidth}{!}{
% \begin{tabular}{l l cccc cc cc}
% \toprule
% \multirow{2}{*}{\textbf{Method}} & \multirow{2}{*}{\textbf{Opt.}} &
% \multicolumn{4}{c}{\textbf{CL Performance} ($\uparrow$)} &
% \multicolumn{2}{c}{\textbf{Weight Flatness} ($\downarrow$)} &
% \multicolumn{2}{c}{\textbf{Feature Robustness}} \\
% \cmidrule(lr){3-6}\cmidrule(lr){7-8}\cmidrule(lr){9-10}
% & &
% $\operatorname{ACC}_{20}^{\mathrm{cls}}$ &
% $\operatorname{AAA}^{\mathrm{cls}}$ &
% $\operatorname{ACC}_{20}^{\mathrm{nme}}$ &
% $\operatorname{AAA}^{\mathrm{nme}}$ &
% Sh$^{(0)}(\rho)$ &
% $\mathrm{tr}(F)$ &
% L2 drift ($\downarrow$) &
% $\mathrm{tr}(E)$ \\
% \midrule
% \multirow{1}{*}{\textbf{PerTaskLP}}
% & SGD
% & $55.22^{\pm 0.31}$ & $59.36^{\pm 0.19}$
% & $62.96^{\pm 0.12}$ & $66.58^{\pm 0.37}$
% & 0.072 & 1015.07
% & 0.00 & 0.82 \\
% \midrule
% \multirow{3}{*}{\textbf{SeqFT}}
% & SGD
% & $50.72^{\pm 2.20}$ & $59.89^{\pm 1.32}$
% & $69.49^{\pm 0.16}$ & $75.10^{\pm 0.57}$
% & 0.582 & 29703.66
% & 29.57 & 0.58 \\
% & SAM
% & $56.07^{\pm 2.60}$ & $66.47^{\pm 1.53}$
% & $72.36^{\pm 2.46}$ & $78.31^{\pm 0.50}$
% & 0.259 & 5881.32
% & 17.04 & 0.74 \\
% & GAM
% & $59.78^{\pm 0.74}$ & $67.31^{\pm 1.18}$
% & $73.74^{\pm 0.88}$ & $78.32^{\pm 0.94}$
% & 0.177 & 3273.08
% & 13.80 & 0.78 \\
% \midrule
% \multirow{3}{*}{\textbf{SeqLoRA}}
% & SGD
% & $63.14^{\pm 2.17}$ & $69.07^{\pm 1.33}$
% & $72.82^{\pm 1.27}$ & $76.62^{\pm 0.68}$
% & 0.142 & 1401.15
% & 16.58 & 0.77 \\
% & SAM
% & $67.04^{\pm 0.71}$ & $71.42^{\pm 0.71}$
% & $74.53^{\pm 0.60}$ & $77.62^{\pm 0.43}$
% & 0.120 & 821.45
% & 14.82 & 0.81 \\
% & GAM
% & $73.31^{\pm 0.52}$ & $76.80^{\pm 0.56}$
% & $75.74^{\pm 0.15}$ & $78.86^{\pm 0.53}$
% & 0.060 & 240.91
% & 11.28 & 0.78 \\
% \midrule
% \multirow{3}{*}{\textbf{IncLoRA}}
% & SGD
% & $65.77^{\pm 2.84}$ & $71.29^{\pm 0.78}$
% & $75.41^{\pm 0.26}$ & $78.37^{\pm 0.26}$
% & 0.076 & 836.91
% & 16.87 & 0.67 \\
% & SAM
% & $68.86^{\pm 0.92}$ & $72.59^{\pm 0.39}$
% & $75.95^{\pm 0.07}$ & $78.68^{\pm 0.25}$
% & 0.053 & 526.49
% & 15.14 & 0.70 \\
% & GAM
% & $72.86^{\pm 0.23}$ & $76.52^{\pm 0.48}$
% & $76.72^{\pm 0.54}$ & $79.42^{\pm 0.40}$
% & 0.035 & 192.25
% & 12.15 & 0.73 \\
% \midrule
% \multirow{3}{*}{\textbf{OLoRA}}
% & SGD
% & $66.64^{\pm 0.44}$ & $71.41^{\pm 0.47}$
% & $75.10^{\pm 0.41}$ & $78.65^{\pm 0.67}$
% & 0.078 & 831.23
% & 17.23 & 0.68 \\
% & SAM
% & $68.53^{\pm 0.78}$ & $72.75^{\pm 0.76}$
% & $75.73^{\pm 0.32}$ & $78.70^{\pm 0.67}$
% & 0.066 & 628.12
% & 15.17 & 0.70 \\
% & GAM
% & $72.44^{\pm 0.36}$ & $76.05^{\pm 0.52}$
% & $76.56^{\pm 0.53}$ & $79.16^{\pm 0.43}$
% & 0.036 & 198.30
% & 12.16 & 0.74 \\
% % \midrule
% % \multirow{1}{*}{\textbf{l2p}}
% % & Adam
% % & $71.35^{\pm 0.05}$ & $74.84^{\pm 0.73}$
% % & -- & --
% % & -- & --
% % & -- & -- \\
% % \midrule
% % \multirow{1}{*}{\textbf{dualprompt}}
% % & Adam
% % & $71.59^{\pm 0.68}$ & $74.80^{\pm 1.18}$
% % & -- & --
% % & -- & --
% % & -- & -- \\
% \bottomrule
% \end{tabular}
% }
% \end{table*}

\subsection{Models and Flatness Optimizers}

All experiments use a ViT-B/16 backbone~\cite{dosovitskiy2020image} pre-trained on ImageNet-21K. To probe the intrinsic representation capacity of the backbone, we include a PerTask Linear Probe baseline that trains an independent linear classifier per task. For PECL, we attach LoRA adapters to each transformer block and study three adapter  representative organizations: SeqLoRA, IncLoRA, and OLoRA.
% SeqLoRA uses a single adapter that is updated across all tasks. IncLoRA incrementally adds task-specific adapters on top of a shared backbone. OLoRA augments this incremental design with orthogonality-promoting constraints. 
All methods share the same backbone and training budget so that differences can be attributed to adapter geometry and flatness optimization.

To promote flatness we integrate several sharpness-based methods. SAM~\cite{foretSharpnessAwareMinimizationEfficiently2021} minimizes a worst-case loss in a small neighborhood around the current parameters.. RWP~\cite{pmlr-v151-bisla22a}, GAM~\cite{zhangGradientNormAware2023}, and C-FLAT~\cite{NEURIPS2024_C-Flat} instantiate different relaxations of sharpness and gradient norm control. In the PECL setting these optimizers act only on the trainable adapter parameters unless explicitly stated. For comparison with global perturbations we additionally consider a Flat-LoRA style variant that applies RWP to all model parameters instead of restricting updates to the adapters.

\subsection{Metrics}

We adopt standard continual-learning metrics based on the task accuracy matrix $a_{tj}$, where $a_{tj}$ denotes the accuracy on task $j$ after learning task $t$. The accuracy at step $t$ is
$
\operatorname{ACC}_t = \frac{1}{t} \sum_{j=1}^{t} a_{tj},
$
and we report final accuracy $\operatorname{ACC}_T^{\mathrm{cls}}$, time-averaged accuracy $\operatorname{AAA}^{\mathrm{cls}} = \frac{1}{T}\sum_{t=1}^{T} \operatorname{ACC}_t.$
% and backward transfer $\operatorname{BWT}^{\mathrm{cls}}$.  , \operatorname{BWT}^{\mathrm{nme}}
In addition, we compute the corresponding Nearest-Class-Mean Classifier (NCM) results
% Nearest Mean of Exemplars (NME) metrics 
$(\operatorname{ACC}_T^{\mathrm{ncm}}, \operatorname{AAA}^{\mathrm{ncm}})$ using class prototypes in the feature space, which isolates the quality and stability of representations from the classifier head.

To characterize flatness and feature stability, we approximate the maximum eigenvalue and the trace of the empirical Fisher information using Lanczos methods~\cite{ghorbaniInvestigationNeuralNet2019,foretSharpnessAwareMinimizationEfficiently2021}.  
% use Lanczos methods~\cite{ghorbaniInvestigationNeuralNet2019,foretSharpnessAwareMinimizationEfficiently2021} to compute the top eigenvalues of Hessian and Hutchinson’s method [4, 5, 65] to compute the Hessian trace.
We further measure the $\ell_2$ drift of class prototypes across tasks and the spectrum of the empirical feature matrix. 
% Together, these quantities provide complementary views on sharpness and feature stability over the task sequence.
% We use power iteration [66] to compute the top eigenvalues of Hessian and Hutchinson’s method [4, 5, 65] to compute the Hessian trace.


\section{Results and Analysis}
\label{sec:results}



% \begin{table*}[htpb]
% \centering
% \caption{
% Class incremental performance, weight-space flatness and feature-space robustness on ImageNet-R ($T=20$, $R=16$). 
% }
% \label{tab:table_all}
% \setlength{\tabcolsep}{4pt}
% \renewcommand{\arraystretch}{0.95}
% \resizebox{\linewidth}{!}{
% \begin{tabular}{l l cccc cc cc}
% \toprule
% \multirow{2}{*}{\textbf{Method}} & \multirow{2}{*}{\textbf{Opt.}} &
% \multicolumn{4}{c}{\textbf{CL Performance} ($\uparrow$)} &
% \multicolumn{2}{c}{\textbf{Weight Flatness} ($\downarrow$)} &
% \multicolumn{2}{c}{\textbf{Feature Robustness}} \\
% \cmidrule(lr){3-6}\cmidrule(lr){7-8}\cmidrule(lr){9-10}
% & &
% $\operatorname{ACC}_{20}^{\mathrm{cls}}$ &
% $\operatorname{AAA}^{\mathrm{cls}}$ &
% $\operatorname{ACC}_{20}^{\mathrm{nme}}$ &
% $\operatorname{AAA}^{\mathrm{nme}}$ &
% Sh$^{(0)}(\rho)$ &
% $\mathrm{tr}(F)$ &
% L2 drift ($\downarrow$) &
% $\mathrm{tr}(E)$ \\
% \midrule
% \multirow{1}{*}{\textbf{PerTaskLP}}
% & SGD
% & $55.22^{\pm 0.31}$ & $59.36^{\pm 0.19}$
% & $62.96^{\pm 0.12}$ & $66.58^{\pm 0.37}$
% & 0.072 & 1015.07
% & 0.00 & 0.82 \\
% \midrule
% \multirow{3}{*}{\textbf{SeqFT}}
% & SGD
% & $50.72^{\pm 2.20}$ & $59.89^{\pm 1.32}$
% & $69.49^{\pm 0.16}$ & $75.10^{\pm 0.57}$
% & 0.582 & 29703.66
% & 29.57 & 0.58 \\
% & SAM
% & $56.07^{\pm 2.60}$ & $66.47^{\pm 1.53}$
% & $72.36^{\pm 2.46}$ & $78.31^{\pm 0.50}$
% & 0.259 & 5881.32
% & 17.04 & 0.74 \\
% & GAM
% & $59.78^{\pm 0.74}$ & $67.31^{\pm 1.18}$
% & $73.74^{\pm 0.88}$ & $78.32^{\pm 0.94}$
% & 0.177 & 3273.08
% & 13.80 & 0.78 \\
% \midrule
% \multirow{3}{*}{\textbf{SeqLoRA}}
% & SGD
% & $63.14^{\pm 2.17}$ & $69.07^{\pm 1.33}$
% & $72.82^{\pm 1.27}$ & $76.62^{\pm 0.68}$
% & 0.142 & 1401.15
% & 16.58 & 0.77 \\
% & SAM
% & $67.04^{\pm 0.71}$ & $71.42^{\pm 0.71}$
% & $74.53^{\pm 0.60}$ & $77.62^{\pm 0.43}$
% & 0.120 & 821.45
% & 14.82 & 0.81 \\
% & GAM
% & $73.31^{\pm 0.52}$ & $76.80^{\pm 0.56}$
% & $75.74^{\pm 0.15}$ & $78.86^{\pm 0.53}$
% & 0.060 & 240.91
% & 11.28 & 0.78 \\
% \midrule
% \multirow{3}{*}{\textbf{IncLoRA}}
% & SGD
% & $65.77^{\pm 2.84}$ & $71.29^{\pm 0.78}$
% & $75.41^{\pm 0.26}$ & $78.37^{\pm 0.26}$
% & 0.076 & 836.91
% & 16.87 & 0.67 \\
% & SAM
% & $68.86^{\pm 0.92}$ & $72.59^{\pm 0.39}$
% & $75.95^{\pm 0.07}$ & $78.68^{\pm 0.25}$
% & 0.053 & 526.49
% & 15.14 & 0.70 \\
% & GAM
% & $72.86^{\pm 0.23}$ & $76.52^{\pm 0.48}$
% & $76.72^{\pm 0.54}$ & $79.42^{\pm 0.40}$
% & 0.035 & 192.25
% & 12.15 & 0.73 \\
% \midrule
% \multirow{3}{*}{\textbf{OLoRA}}
% & SGD
% & $66.64^{\pm 0.44}$ & $71.41^{\pm 0.47}$
% & $75.10^{\pm 0.41}$ & $78.65^{\pm 0.67}$
% & 0.078 & 831.23
% & 17.23 & 0.68 \\
% & SAM
% & $68.53^{\pm 0.78}$ & $72.75^{\pm 0.76}$
% & $75.73^{\pm 0.32}$ & $78.70^{\pm 0.67}$
% & 0.066 & 628.12
% & 15.17 & 0.70 \\
% & GAM
% & $72.44^{\pm 0.36}$ & $76.05^{\pm 0.52}$
% & $76.56^{\pm 0.53}$ & $79.16^{\pm 0.43}$
% & 0.036 & 198.30
% & 12.16 & 0.74 \\
% % \midrule
% % \multirow{1}{*}{\textbf{l2p}}
% % & Adam
% % & $71.35^{\pm 0.05}$ & $74.84^{\pm 0.73}$
% % & -- & --
% % & -- & --
% % & -- & -- \\
% % \midrule
% % \multirow{1}{*}{\textbf{dualprompt}}
% % & Adam
% % & $71.59^{\pm 0.68}$ & $74.80^{\pm 1.18}$
% % & -- & --
% % & -- & --
% % & -- & -- \\
% \bottomrule
% \end{tabular}
% }
% \end{table*}




\subsection{Main results on ImageNet-R}
\label{sec:results-main-inr}

\begin{figure*}[thbp]
    \centering
    % ImageNet-C: cls / nme
    % \begin{minipage}{0.495\linewidth}
    %     \centering
    %     \includegraphics[width=0.85\linewidth]{figures_TRML/AAA_curve_Imagenet-C_summary_inc10_rank16_ponit2.pdf}
    %     % \subcaption{ImageNet-C: $\operatorname{ACC}_t^{\mathrm{cls}}$}
    %     \label{fig:aaa_imagenet_c}
    % \end{minipage}\hfill
    \begin{minipage}{0.495\linewidth}
        \centering
        \includegraphics[width=0.85\linewidth]{figures_TRML/NME_curve_Imagenet-C_summary_inc10_rank16_point2.pdf}
        % \subcaption{ImageNet-C: $\operatorname{ACC}_t^{\mathrm{nme}}$}
        \label{fig:nme_imagenet_c}
    \end{minipage}\hfill
    % ImageNet-P: cls / nme
    % \begin{minipage}{0.495\linewidth}
    %     \centering
    %     \includegraphics[width=0.85\linewidth]{figures_TRML/AAA_curve_Imagenet-P_summary_inc10_rank16_point2.pdf}
    %     % \subcaption{ImageNet-P: $\operatorname{ACC}_t^{\mathrm{cls}}$}
    %     \label{fig:aaa_imagenet_p}
    % \end{minipage}\hfill
    \begin{minipage}{0.495\linewidth}
        \centering
        \includegraphics[width=0.85\linewidth]{figures_TRML/NME_curve_Imagenet-P_summary_inc10_rank16_point2.pdf}
        % \subcaption{ImageNet-P: $\operatorname{ACC}_t^{\mathrm{nme}}$}
        \label{fig:nme_imagenet_p}
    \end{minipage}
    \caption{ 
        The full trajectories
of$\operatorname{ACC}_t^{\mathrm{ncm}}$ on ImageNet-C and ImageNet-P. 
        % The four panels report performance after each task for the same class incremental methods and optimization configurations as in Figure~\ref{fig:combined_curves_imageR}. 
         SAM and GAM consistently outperform SGD on both benchmark, which indicates that promoting flat minima in the adapter improves robustness of the learned representations.
         %   and $\operatorname{ACC}_t^{\mathrm{cls}}$
    }
    \label{fig:robustness_curves_comparison}
\end{figure*}

 
% Table~\ref{tab:table_all} summarises final accuracy $\operatorname{ACC}_{20}^{\mathrm{cls}}$, averaged accuracy $\operatorname{AAA}^{\mathrm{cls}}$, with the corresponding NME metrics together with representative sharpness proxies in the parameter spac and feature space diagnostics. Figure~\ref{fig:combined_curves_imageR} reports classifier head accuracy $\operatorname{ACC}_t^{\mathrm{cls}}$ and NME based accuracy $\operatorname{ACC}_t^{\mathrm{nme}}$ over the task sequence on ImageNet-R.

Table~\ref{tab:table_all} reports, on ImageNet-R, the final and time-averaged accuracies for both classifier heads and NCM evaluation, together with weight-space flatness proxies (adversarial sharpness, largest eigenvalue and trace of the empirical Fisher) and feature-space diagnostics (prototype drift and feature-matrix trace). 
% Figure~\ref{fig:combined_curves_imageR} complements this summary with the full trajectories of $\operatorname{ACC}_t^{\mathrm{cls}}$ and $\operatorname{ACC}_t^{\mathrm{nme}}$ over the task sequence.


Across all three LoRA organizations, applying SAM or GAM to the adapter parameters consistently improves class-incremental performance compared with SGD under the same training budget. Both final and time-averaged accuracies increase for the classifier head and for NCM evaluation, that is, $\operatorname{ACC}_{20}^{\mathrm{cls}}$, $\operatorname{AAA}^{\mathrm{cls}}$, $\operatorname{ACC}_{20}^{\mathrm{ncm}}$ and $\operatorname{AAA}^{\mathrm{ncm}}$ all improve, with GAM typically achieving the strongest gains. 
% The trajectories in Figure~\ref{fig:combined_curves_imageR} show that this advantage is in all process: SAM and GAM maintain higher accuracy throughout the sequence. The  alignment between the classifier-head and NME curves indicates that these gains reflect more stable shared representations.
These performance gains are tightly aligned with our flatness and feature metrics. For each LoRA organization, moving from SGD to SAM and then to GAM monotonically decreases adversarial sharpness $\mathrm{Sh}^{(0)}(\rho)$ as well as the largest eigenvalue and trace of the empirical Fisher, while simultaneously reducing prototype drift, and yielding a more isotropic feature matrix. 

Taken together, these trends support the central prediction of our PAC-Bayesian analysis that reducing expected sharpness within the adapter subspace is sufficient to tighten the generalization guarantee and improve the stability–plasticity trade-off in PECL.

% \begin{figure}[hpbt]
%     \centering
%     \begin{minipage}{0.49\linewidth}
%         \centering
%         \includegraphics[width=\linewidth]{figures_TRML/AAA_curve_Imagenet-R_summary_inc10_rank16—point2.pdf}
%         % \subcaption{Classifier head accuracy $\operatorname{ACC}_t^{\mathrm{cls}}$ on ImageNet-R}
%         \label{fig:curves_imageR_1_AAA}
%     \end{minipage}
%      \begin{minipage}{0.49\linewidth}
%         \centering
%         \includegraphics[width=\linewidth]{figures_TRML/NME_curve_Imagenet-R_summary_inc10_rank16_point2.pdf}
%         % \subcaption{Prototype based accuracy $\operatorname{ACC}_t^{\mathrm{nme}}$ on ImageNet-R}
%         \label{fig:curves_imageR_2_NME}
%     \end{minipage}
%     \caption{
%     The full trajectories
% of $\operatorname{ACC}_t^{\mathrm{cls}}$ and $\operatorname{ACC}_t^{\mathrm{nme}}$ on ImageNet-R. 
%     % Summary curves of classifier head accuracy $\operatorname{ACC}_t^{\mathrm{cls}}$ and NME based accuracy $\operatorname{ACC}_t^{\mathrm{nme}}$ on ImageNet-R. 
%     % For all three LoRA organizations, flatness promoting optimizers such as SAM and GAM consistently outperform SGD, and their advantage becomes more pronounced on later tasks, which indicates more stable representations and reduced forgetting.
%     }
%     \label{fig:combined_curves_imageR}
% \end{figure}








% tables \ref{table:q1_ci_with_weight_flatness_reorg} and \ref{tab:q1_lp_feat_reorg} remain as in Section~\ref{sec: flat minima}. We only change how they are discussed here.



\subsection{Robustness on ImageNet-C and ImageNet-P}
\label{sec:results-robustness}

We next evaluate robustness under corruptions and perturbations on Tiny ImageNet-C and Tiny ImageNet-P. Fig.~\ref{fig:robustness_curves_comparison} shows the trajectories of $\operatorname{ACC}_t^{\mathrm{ncm}}$ over the task sequence on these benchmarks. 

% $\operatorname{ACC}_t^{\mathrm{cls}}$ and 

Across a wide range of corruption types and perturbation sequences, sharpness-based LoRA training consistently improves performance: SAM and GAM achieve higher $\operatorname{ACC}_t^{\mathrm{cls}}$ and $\operatorname{ACC}_t^{\mathrm{ncm}}$ than SGD. This pattern indicates that promoting flatness in the LoRA subspace does not trade clean accuracy for robustness but instead enhances generalization in a way that transfers to distribution shift. SeqLoRA becomes the most robust configuration, while IncLoRA and OLoRA suffer larger accuracy drops. This is consistent with the dependence of the effective sharpness term on the spectrum of $F_W$ restricted to the adapter subspace. Larger update subspaces admit more directions along which curvature and shift induced drift can accumulate, whereas the compact SeqLoRA subspace constrains both $\Sigma_t$ and $F_\Delta$ to a smaller region and therefore exposes fewer unstable directions.


% \begin{figure}[htbp]
%     \centering
%     % ImageNet-C: cls / nme
%     \begin{minipage}{0.48\linewidth}
%         \centering
%         \includegraphics[width=\linewidth]{figures_TRML/AAA_curve_Imagenet-C_summary_inc10_rank16_ponit2.pdf}
%         % \subcaption{ImageNet-C: $\operatorname{ACC}_t^{\mathrm{cls}}$}
%         \label{fig:aaa_imagenet_c}
%     \end{minipage}\hfill
%     % ImageNet-P: cls / nme
%     \begin{minipage}{0.48\linewidth}
%         \centering
%         \includegraphics[width=\linewidth]{figures_TRML/AAA_curve_Imagenet-P_summary_inc10_rank16_point2.pdf}
%         % \subcaption{ImageNet-P: $\operatorname{ACC}_t^{\mathrm{cls}}$}
%         \label{fig:aaa_imagenet_p}
%     \end{minipage}\hfill
    
%     \caption{
%         The NME-based ($\operatorname{ACC}_t^{\mathrm{nme}}$) on ImageNet-C and ImageNet-P. 
%         % The four panels report performance after each task for the same class incremental methods and optimization configurations as in Figure~\ref{fig:combined_curves_imageR}. 
%          SAM and GAM consistently outperform SGD on both corruption robustness (ImageNet-C) and perturbation robustness (ImageNet-P), which indicates that promoting flat minima in the adapter improves robustness of the learned representations.
%     }
%     \label{fig:robustness_curves_comparison}
% \end{figure}


% We next evaluate robustness under corruptions and perturbations on Tiny ImageNet-C and Tiny ImageNet-P.
%  Figure~\ref{fig:robustness_curves_comparison} complements this summary with the full trajectories of $\operatorname{ACC}_t^{\mathrm{cls}}$ and $\operatorname{ACC}_t^{\mathrm{nme}}$ over the task sequence.
% % For each LoRA organization and optimizer, we measure performance using $\operatorname{ACC}_t^{\mathrm{cls}}$ and $\operatorname{ACC}_t^{\mathrm{nme}}$.

% Across a wide range of corruption types and perturbation sequences, flatness oriented LoRA training consistently improves robustness. 
% SAM and GAM achieve higher corrupted $\operatorname{ACC}_t^{\mathrm{cls}}$ and $\operatorname{ACC}_t^{\mathrm{nme}}$ than SGD, and the ordering observed on clean ImageNet-R is preserved on ImageNet-C and ImageNet-P. 
% The gains under corruption are strongly correlated with the gains on clean data: configurations that achieve higher clean AAA also yield higher corrupted AAA. 
% This pattern indicates that LoRA subspace flatness does not trade clean accuracy for robustness. 
% Instead, it leads to a global enhancement of generalization that extends to distribution shifts, which is consistent with the PAC Bayesian view that smaller expected sharpness tightens the error bound uniformly over the data distribution.



\subsection{Local versus global perturbations}
\label{sec:results-local-global}



We examine whether flatness should be promoted locally in the adapter subspace or globally in the full parameter space in PECL. 
We compare the Flat-LoRA configuration, which applies RWP to all model parameters, with a PECL compatible
variant that restricts RWP to the LoRA branch only. The experiment is repeated for the three LoRA organizations and different perturbation strengths and results are shown in Table\ref{tab:lora——or_full}.


\begin{table}[hpb]
\centering
\caption{$\operatorname{AAA}^{\mathrm{cls}}$ and $\operatorname{AAA}^{\mathrm{ncm}}$ on ImageNet-R under global (full model) versus subspace (LoRA only).}
\label{tab:lora——or_full}
\setlength{\tabcolsep}{3pt}
\renewcommand{\arraystretch}{1}
\resizebox{1\linewidth}{!}{
\begin{tabular}{c c cc cc cc}
\toprule
\multicolumn{2}{c}{\textbf{Opt.}} &
\multicolumn{2}{c}{\textbf{SeqLoRA}} &
\multicolumn{2}{c}{\textbf{IncLoRA}} &
\multicolumn{2}{c}{\textbf{OLoRA}} \\
\cmidrule(lr){1-2} \cmidrule(lr){3-4} \cmidrule(lr){5-6} \cmidrule(lr){7-8}
\textbf{Strength} & \textbf{Range} &
$\operatorname{AAA}^{\mathrm{cls}}$ & $\operatorname{AAA}^{\mathrm{ncm}}$ &
$\operatorname{AAA}^{\mathrm{cls}}$ & $\operatorname{AAA}^{\mathrm{ncm}}$ &
$\operatorname{AAA}^{\mathrm{cls}}$ & $\operatorname{AAA}^{\mathrm{ncm}}$ \\
\midrule
\multirow{1}{*}{$0$} &Base &
69.07 & 76.62 & 71.29 & 78.29 & 71.41 & 78.47 \\
\midrule
% 组 1：\rho_1
\multirow{2}{*}{$0.05$}
&Full 
&-3.5  & -1.54 & -2    & -0.38 & -2.62 & -0.8  \\ 
& LoRA &0.27  & 0.02  & 0.79  & 0.51  & 0.12  & 0.16  \\ 
\midrule
% 组 2：\rho_2
\multirow{2}{*}{$0.10$} & Full &-4.82 & -1.29 & -4.74 & -0.74 & -5.53 & -1.18 \\
& LoRA &0.7   & 0.29  & 0.86  & 0.53  & 0.28  & 0.12  \\
\midrule
% 组 3：\rho_3
\multirow{2}{*}{$0.15$} & Full &-6.72 & -1.69 & -6.84 & -0.98 & -6.49 & -1.29 \\
& LoRA &0.41  & 0.08  & 0.77  & 0.53  & 0.27  & 0.11 \\
\bottomrule
\end{tabular}}
\end{table}


Across all settings, global perturbations are uniformly harmful in the PECL regime. 
Applying RWP to the full model substantially degrades both $\operatorname{AAA}^{\mathrm{cls}}$ and $\operatorname{AAA}^{\mathrm{ncm}}$ on ImageNet-R for all LoRA organizations and perturbation ranges, often performing worse than the plain SGD baseline. 
In contrast, restricting RWP to the LoRA subspace yields small but consistent improvements over SGD and matches or slightly trails the gains obtained by SAM and GAM.

These findings are aligned with the structural assumptions in our theory. 
In PECL, the posterior covariance is confined to the adapter subspace $W_\Delta$, while the backbone and past adapters are treated as deterministic and are not supposed to carry posterior uncertainty. 
Global perturbations violate this assumption, inject noise into parameters that encode past knowledge, and empirically disrupt the stability of previously learned tasks. 
By comparison, local perturbations confined to $W_\Delta$ respect the PECL structure: they reduce the subspace expected sharpness that appears in the PAC Bayesian bound without disturbing the frozen components. 
This explains why Flat LoRA style full perturbations are not suitable in the PECL.




\subsection{Comparing flatness optimizers}
\label{sec:results-optimizers}



\begin{table}[thpb]
\centering
\caption{The $\operatorname{AAA}^{\mathrm{cls}}$ and $\operatorname{AAA}^{\mathrm{ncm}}$ performance across LoRA organizations under different flatness optimizers. 
% Rows below SGD report the change relative to the SGD baseline.
}
\label{tab:aaa_nme_lora_org}
% \setlength{\tabcolsep}{5pt}
% \resizebox{\linewidth}{!}
\setlength{\tabcolsep}{3pt}
\renewcommand{\arraystretch}{1} % 可选：稍微压缩行距
\resizebox{1\linewidth}{!}{
\begin{tabular}{l cc cc cc}
\toprule
\multirow{3}{*}{\textbf{Opt.}} &
\multicolumn{2}{c}{\textbf{SeqLoRA}} &
\multicolumn{2}{c}{\textbf{IncLoRA}} &
\multicolumn{2}{c}{\textbf{OLoRA}} \\
\cmidrule(lr){2-3} \cmidrule(lr){4-5} \cmidrule(lr){6-7}
% & $\overline{\text{ACC}}_{20}$ (CL) & $\overline{\text{ACC}}_{20}$ (NME) & $\overline{\text{ACC}}_{20}$ (CL) & $\overline{\text{ACC}}_{20}$ (NME) & $\overline{\text{ACC}}_{20}$ (CL) & $\overline{\text{ACC}}_{20}$ (NME) \\
% \cmidrule(lr)
% &  CL &  NME & CL &  NME &  CL &  NME \\
% & $\overline{\text{ACC}}_{20}$ & $\overline{\text{ACC}}_{20}$  & $\overline{\text{ACC}}_{20}$  & $\overline{\text{ACC}}_{20}$  & $\overline{\text{ACC}}_{20}$  & $\overline{\text{ACC}}_{20}$  \\
& $\operatorname{AAA}^{\mathrm{cls}}$ &  $\operatorname{AAA}^{\mathrm{ncm}}$ & $\operatorname{AAA}^{\mathrm{cls}}$ & $\operatorname{AAA}^{\mathrm{ncm}}$ & $\operatorname{AAA}^{\mathrm{cls}}$ & $\operatorname{AAA}^{\mathrm{ncm}}$ \\
\midrule
% $\overline{\text{ACC}}_{20}$
SGD  & 69.07 & 76.62 & 71.29 & 78.29 & 71.41 & 78.47 \\
\midrule
$\Delta$ RWP   &0.70  & 0.29  & 0.86  & 0.53  & 0.28  & 0.12   \\
$\Delta$ SAM   & 2.35  & 1.00  & 1.29  & 0.40  & 1.08  & 0.22  \\
$\Delta$ CFLAT & 2.35  & 1.04  & 2.67  & 1.26  & 1.07  & 0.12  \\
$\Delta$ GAM  
% & 7.73  & 2.24  & 5.23  & 1.13  & 4.64  & 0.68
& \textbf{7.73}  & \textbf{2.24}  & \textbf{5.23}  & \textbf{1.13}  & \textbf{4.64}  & \textbf{0.68}


\\
\bottomrule
\end{tabular}}
\end{table}

To isolate the effect of different flatness optimizers, we fix the LoRA organization and hyperparameters and compare RWP, SAM, C-FLAT, and GAM applied only to the adapter parameters, with SGD as a baseline. 


% The comparison uses the same ImageNet-R PECL setting as in Section~\ref{sec:results-main-inr} and reports $\operatorname{AAA}^{\mathrm{cls}}$, $\operatorname{AAA}^{\mathrm{nme}}$, and a subset of sharpness metrics.

The results in Table\ref{tab:aaa_nme_lora_org} exhibit a clear ordering. 
RWP only slightly above  SGD. 
SAM and C FLAT achieve intermediate gains. 
GAM consistently  and the highest $\operatorname{AAA}^{\mathrm{cls}}$ and $\operatorname{AAA}^{\mathrm{ncm}}$. 
The performance is also aligned the performance in Fig\ref{fig:robustness_curves_comparison} with sharpness result in Table\ref{tab:table_all}.
This ordering reveals a parallel hierarchy on the geometric side: optimisers that more aggressively suppress subspace sharpness also achieve better continual learning performance. In other words, the trend appears simultaneously in the flatness diagnostics and in the CI metrics. This monotone relationship is precisely what Theorem~\ref{thm:pecl-bound} predicts: for a fixed LoRA geometry and training budget, methods that reduce the adapter-subspace sharpness tighten the empirical term of the PAC–Bayesian bound and allow the posterior and hyper-posterior KL terms to remain moderate, which manifests empirically as higher $\operatorname{AAA}^{\mathrm{cls}}$ and $\operatorname{AAA}^{\mathrm{ncm}}$.

\subsection{Result on NLP Benchmark}


We further conduct continual task learning experiments on two standard NLP benchmarks under multiple task orders, and report the final accuracy $ACC_T$. Table~\ref{tab:nlp_accT_orders} shows that SAM consistently improves $ACC_T$ over the SGD baselines across most configurations. The gains are more pronounced in the long-sequence regime, suggesting that promoting flatness within the LoRA subspace remains effective when inter-task interference accumulates. Although a few order-specific deviations are observed, the overall trend indicates that SAM is a generally beneficial for LoRA-based CL methods in NLP.


\begin{table}[thb]
\centering
\caption{Summary of the results on two standard CL benchmarks with T5-large model. $ACC_T$ is reported.} % 标题不改：如你原本有 caption 就把它放这里
\label{tab:nlp_accT_orders} % 如你原本有 label 就沿用
% 列间距和行高压缩
\setlength{\tabcolsep}{3pt}
\renewcommand{\arraystretch}{0.9}
\small
\resizebox{0.98\linewidth}{!}{
\begin{tabular}{l cccc cccc}
\toprule
\multirow{2}{*}{} &
\multicolumn{4}{c}{\textbf{Standard CL Benchmark}} &
\multicolumn{4}{c}{\textbf{Large Number of Tasks}} \\
\cmidrule(lr){2-5}\cmidrule(lr){6-9}
\textbf{Method} &
\textbf{Order1} & \textbf{Order2} & \textbf{Order3} & \textbf{AVG} &
\textbf{Order4} & \textbf{Order5} & \textbf{Order6} & \textbf{AVG} \\
\midrule
SeqLoRA     & 61.66          & 59.99          & 72.22          & 64.62          & 55.85          & 40.65          & 27.05          & 41.18         \\
SeqLoRA-SAM & \textbf{63.33} & \textbf{63.16} & \textbf{72.95} & \textbf{66.48} & \textbf{56.20} & \textbf{56.25} & \textbf{12.77} & \textbf{41.74}   \\
\midrule
IncLoRA     &\textbf{71.18}          & \textbf{71.56}          & 72.87          & 71.87          & 58.13          & 56.42          & 39.26          & 51.27            \\
IncLoRA-SAM & 70.20          & 71.18          & \textbf{75.78} & \textbf{72.39}          & \textbf{67.5}  & \textbf{56.54}          & \textbf{66.38}          & \textbf{63.47} \\
\midrule
OLoRA       & 76.41          & 76.93          & 73.69          & 75.68          & 66.73          & 68.62          &\textbf{71.51}          & 68.95         \\
OLoRA-SAM   & \textbf{78.31} & \textbf{77.96} & \textbf{75.95} & \textbf{77.41} & \textbf{71.08} & \textbf{69.05} & {70.08} & \textbf{70.07}   \\
\bottomrule
\end{tabular}
}
\end{table}


\subsection{Ablation on task length and adapter rank}
\label{sec:results-ablation}


We further vary the task sequence length and the adapter rank on ImageNet-R and report $\operatorname{AAA}^{\mathrm{ncm}}$ in Fig.~\ref{fig:ablation}. Across all values of $T$ and $r$, the relative behavior observed remains stable: flatness aware optimizers consistently outperform or SGD counterparts for SeqLoRA, IncLoRA, and OLoRA.

% Finally, we study the robustness of the above conclusions with respect to the task sequence length $T$ and the LoRA rank $r$. 
% We repeat the ImageNet-R experiments for shorter and longer sequences and for several ranks $r$ while keeping all other hyperparameters fixed. 
% The detailed $\operatorname{ACC}_t^{\mathrm{cls}}$ and $\operatorname{ACC}_t^{\mathrm{nme}}$ curves are reported in Appendix~\ref{app:ablation-length-rank}.

\begin{figure}[htbp]
    \centering
    \begin{minipage}{0.495\linewidth}
        \centering
        \includegraphics[width=\linewidth]{figures_TRML/length_NME_T_avg_by_method_opt_imagenetr.pdf}
        % \subcaption{Different Length}
        % \label{fig:curves_imageR_length}
    \end{minipage}%
    \begin{minipage}{0.495\linewidth}
        \centering
        \includegraphics[width=\linewidth]{figures_TRML/rank_NME_T_avg_by_method_opt_imagenetr.pdf}
        % \subcaption{Different Rank}
        % \label{fig:curves_imageR_rank}
    \end{minipage}
    \caption{
        Effect of task sequence length and adapter rank on NCM average accuracy on ImageNet-R..
    }
    \label{fig:ablation}
\end{figure}

% The qualitative picture remains stable. 
% For all tested values of $T$ and $r$, applying SAM or GAM in the adapter subspace consistently improves $\operatorname{AAA}^{\mathrm{cls}}$ and $\operatorname{AAA}^{\mathrm{nme}}$ over SGD, and the accuracy gap tends to be more pronounced on later tasks.




% Changing $T$ or $r$ affects the absolute level of performance and the magnitude of forgetting, as expected, but does not alter the relative ordering among optimizers. 
% Feature space diagnostics follow the same pattern: when flatness is promoted in $W_\Delta$, prototype drift is reduced and the empirical feature matrix exhibits higher effective rank across all ablation settings. 

% These results indicate that the empirical support for the sharpness aware PAC Bayesian analysis is robust to reasonable variations in task length and adapter capacity, and does not depend on a narrow choice of hyperparameters.



% \section{Discussion, Limitations, and Related Work}


\section{Conclusion, Limitations and Future Work}
This work revisits PECL through the geometry of LoRA adapter subspaces. Within a hierarchical PAC Bayesian framework, we derive a generalization bound in which the continual learning risk is controlled by an empirical term under Gaussian perturbations, the divergence $\mathrm{KL}(Q_\Delta \Vert P_\Delta)$, and the hyper posterior drift across tasks. By restricting updates to the LoRA manifold $W_\Delta$ and working in factor coordinates, the analysis identifies subspace sharpness as the relevant notion of flatness and links it to the spectrum of a LoRA restricted Fisher matrix and to the stability of feature representations.

The framework is deliberately aligned with a PECL regime where the pretrained backbone and previously learned adapters remain frozen. It does not cover partial or periodic finetuning of the backbone. Extending the analysis to such settings is a natural direction for future work and may reveal additional interactions between subspace geometry and global updates. Even within its current scope, the results suggest several practical guidelines for LoRA based PECL: flatness control should primarily act in the trainable adapter subspace, the organisation of adapters can be viewed as a structural prior that shapes the Fisher spectrum and stability plasticity trade offs, and subspace sharpness together with Fisher based diagnostics provides a useful lens for interpreting and comparing PECL methods.

\bibliography{example_paper}
\bibliographystyle{icml2025}

\clearpage      % 换页，避免和正文混排
\onecolumn      % 从这里开始改为单栏
\appendix       % 如果你需要正式标记为附录，可加这一句

\section{Appendix Proof}


\subsection{Proof of Lemma~\ref{lem:kl-decomp-pecl}}
\label{proof:Lemma1}

% \begin{lemma}[KL decomposition with a frozen backbone]
% \label{equation: subspace-full}
% Let $W\in\mathbb R^d$ be fixed and let $P^{\Delta},Q^{\Delta}$ be probability measures on $(\mathcal W_\Delta,\mathcal B(\mathcal W_\Delta))$ with $Q^{\Delta}\ll P^{\Delta}$. Define their pushforwards onto the affine subspace by
% \[
% P^{\mathrm{full}}:=(T_W)_{\text{\#}} P^{\Delta},\qquad
% Q^{\mathrm{full}}:=(T_W)_{\text{\#}} Q^{\Delta}.
% \]
% Then $Q^{\mathrm{full}}\ll P^{\mathrm{full}}$ and
% \[
% \mathrm{KL}\left(Q^{\mathrm{full}}\middle|P^{\mathrm{full}}\right)
% =\mathrm{KL}\left(Q^{\Delta}\middle|P^{\Delta}\right).
% \]
% Equivalently, identifying ${W}\times\mathcal W_\Delta$ with $W+\mathcal W_\Delta$ via $\Phi$, one may write
% $P^{\mathrm{full}}=\Phi_{\text{\#}}(\delta_W\otimes P^{\Delta})$ and
% $Q^{\mathrm{full}}=\Phi_{\text{\#}}(\delta_W\otimes Q^{\Delta})$, and the same equality holds.
% \end{lemma}

\begin{lemma}[KL in PECL is adapter-subspace KL]
Assume $Q_t^\Delta \ll P_t^\Delta$. Then $Q_t^{\mathrm{full}} \ll P_t^{\mathrm{full}}$ and
$$
\mathrm{KL}\bigl(Q_t^{\mathrm{full}} \,\Vert\, P_t^{\mathrm{full}}\bigr)
\;=\;
\mathrm{KL}\bigl(Q_t^\Delta \,\Vert\, P_t^\Delta\bigr).
$$
\end{lemma}

\begin{proof}
  Absolute continuity is preserved by pushforward under measurable maps, hence $Q^{\mathrm{full}}\ll P^{\mathrm{full}}$. Since $T_W$ is a Borel isomorphism, by the Radon-Nikodym chain rule
 \[
 \frac{dQ^{\mathrm{full}}}{dP^{\mathrm{full}}}(T_W\Delta)=\frac{dQ^{\Delta}}{dP^{\Delta}}(\Delta)\quad\text{for }P^{\Delta}\text{-a.e.\ }\Delta.
 \]

\begin{equation*}
     \begin{aligned}
 \mathrm{KL}\left(Q^{\mathrm{full}}\middle|P^{\mathrm{full}}\right)
 &= \int_{W+\mathcal W_\Delta}
 \log\Big(\frac{dQ^{\mathrm{full}}}{dP^{\mathrm{full}}}(y)\Big)
 Q^{\mathrm{full}}(dy)\\
 &= \int_{\mathcal W_\Delta}
 \log\Big(\frac{dQ^{\mathrm{full}}}{dP^{\mathrm{full}}}\big(T_W\Delta\big)\Big)
 Q^{\Delta}(d\Delta) \quad \text{(} Q^{\mathrm{full}}=(T_W)_{\text{\#}}Q^{\Delta}\text{)} 
 \\
 &= \int_{\mathcal W_\Delta}
 \log\Big(\frac{dQ^{\Delta}}{dP^{\Delta}}(\Delta)\Big)
 Q^{\Delta}(d\Delta) \\
 &= \mathrm{KL}\left(Q^{\Delta}\middle|P^{\Delta}\right).
 \end{aligned}
\end{equation*} 
\end{proof}

% \section{ Proof of Lemma~\ref{lemma: subspace-full}}
% \begin{lemma}[Sharpness decomposition in PECL]
% \label{lemma: Sharpness in PECL}
% Let $W_0\in\mathbb R^d$ be fixed and let the trainable update lie in the adapter subspace $\mathcal W_\Delta\subseteq\mathbb R^d$. Consider a Gaussian posterior on $(\mathcal W_\Delta,\mathcal B(\mathcal W_\Delta))$,
%  $
%  Q_t^{\Delta}=\mathcal N(\hat{\Delta W}_t,\Sigma_t),
%  $
%  with mean $\hat{\Delta W}_t\in\mathcal W_\Delta$ and covariance operator $\Sigma_t\succeq 0$ acting on $\mathcal W_\Delta$ (viewed in $\mathbb R^d$, this covariance is supported in $\mathcal W_\Delta$). Define the full-space posterior on ${W_0}\times\mathcal W_\Delta$ by
%  $
%  Q_t^{\mathrm{full}}:=\delta_{W_0}\otimes Q_t^{\Delta},
%  $
% and identify ${W_0}\times\mathcal W_\Delta$ with the affine subspace $W_0+\mathcal W_\Delta$ via $(w,\Delta)\mapsto w+\Delta$. Let $\hat L_{S_t}:\mathbb R^d\to\mathbb R$ be measurable and integrable under $Q_t^{\mathrm{full}}$. Then
%  \[
%  \mathbb{E}_{W\sim Q_t^{\mathrm{full}}}\big[\hat{L}_{S_t}(W)\big]
%  = \mathbb{E}_{\varepsilon\sim\mathcal N(0,\Sigma_t)}
%  \big[\hat{L}_{S_t}(W_0+\hat{\Delta W}_t+\varepsilon)\big]
%  = \hat{L}_{S_t}(W_0+\hat{\Delta W}_t) +\mathrm{E\text{-}Sh}_t^{\Delta}(\hat{\Delta W}_t,\Sigma_t) \]
%    where the subspace expected sharpness is
%    \[
%    \mathrm{E\text{-}Sh}_t^{\Delta}(\hat{\Delta W}_t,\Sigma_t)
%    = \mathbb{E}_{\varepsilon\sim\mathcal N(0,\Sigma_t)}
%    \big[\hat{L}_{S_t}(W_0+\hat{\Delta W}_t+\varepsilon)
%   - \hat{L}_{S_t}(W_0+\hat{\Delta W}_t)\big].
%      \]
% \end{lemma}
% \begin{proof}
%     Sampling $W\sim Q_t^{\mathrm{full}}$ is equivalent to sampling $\varepsilon\sim\mathcal N(0,\Sigma_t)$ on $\mathcal W_\Delta$ and setting $W=W_0+\hat{\Delta W}_t+\varepsilon$. Taking expectations and adding-subtracting $\hat L_{S_t}(W_0+\hat{\Delta W}_t)$ inside the expectation gives the decomposition. 
% \end{proof}

\subsection{Proof of Lemma~\ref{lem:subspace-loss-pecl}}
\label{proof:lemma2}

% \begin{equation*}
% \label{eq:subspace-loss-pecl}
% \begin{aligned}
% &\mathbb E_{W\sim Q_t^{\mathrm{full}}}\bigl[\hat L_{S_t}(W)\bigr] \\
% &= \mathbb{E}_{\varepsilon\sim\mathcal N(0,\Sigma_t)}
%     \big[\hat{L}_{S_t}(W_\mathrm {frozen }+\hat{\Delta W}_t+\varepsilon)\big],\varepsilon\in\mathcal W_\Delta \\
% & =\hat L_{S_t}(W_{\mathrm{frozen}} + \hat{\Delta W}_t)
% +
% \mathrm{E\text{-}Sh}_t^\Delta(\hat{\Delta W}_t,\Sigma_t),\varepsilon\in\mathcal W_\Delta.
% \end{aligned}
% \end{equation*}

\begin{lemma}
    The posterior expected empirical loss decomposes as
\begin{equation*}
    \mathbb{E}_{W\sim Q_t^{\mathrm{full}}}\big[\hat{L}_{S_t}(W)\big]
= \mathbb{E}_{\varepsilon\sim\mathcal N(0,\Sigma_t)}
    \big[\hat{L}_{S_t}(W_\mathrm {frozen }+\hat{\Delta W}_t+\varepsilon)\big]
= \hat{L}_{S_t}(W_\mathrm {frozen }+\hat{\Delta W}_t)
  + \mathrm{E\text{-}Sh}_t^{\Delta}(\hat{\Delta W}_t,\Sigma_t),
\end{equation*}
where the expected sharpness is
$$
\mathrm{E\text{-}Sh}_t^{\Delta}(\hat{\Delta W}_t,\Sigma_t)
:= \mathbb{E}_{\varepsilon\sim\mathcal N(0,\Sigma_t)}
   \big[\hat{L}_{S_t}(W_\mathrm {frozen }+\hat{\Delta W}_t+\varepsilon)
       - \hat{L}_{S_t}(W_\mathrm {frozen }+\hat{\Delta W}_t)\big],\quad
\varepsilon\sim\mathcal N(0,\Sigma_t)\ \text{on } \mathcal W_\Delta$$
\end{lemma}


\begin{proof}
Recall that the adapter posterior is a Gaussian supported on the update subspace,
$
Q_t^\Delta = \mathcal N(\hat{\Delta W}_t,\Sigma_t)
$
on $(\mathcal W_\Delta,\mathcal B(\mathcal W_\Delta))$.
Hence we can represent a draw $\Delta W \sim Q_t^\Delta$ as
$
\Delta W = \hat{\Delta W}_t + \varepsilon
$
with $\varepsilon \sim \mathcal N(0,\Sigma_t)$ on $\mathcal W_\Delta$.

Under the PECL constraint, the full parameter is obtained by translating the adapter update by the frozen part,
$
W = W_{\mathrm{frozen}} + \Delta W.
$
Therefore, for the empirical risk,
\begin{align*}
\mathbb E_{W \sim Q_t^{\mathrm{full}}}\big[\hat L_{S_t}(W)\big]
&= \mathbb E_{\Delta W \sim Q_t^\Delta}\big[\hat L_{S_t}(W_{\mathrm{frozen}}+\Delta W)\big] \\
&= \mathbb E_{\varepsilon \sim \mathcal N(0,\Sigma_t)}
\big[\hat L_{S_t}(W_{\mathrm{frozen}}+\hat{\Delta W}_t+\varepsilon)\big].
\end{align*}
Finally, add and subtract the mean point loss to obtain
\begin{align*}
\mathbb E_{\varepsilon \sim \mathcal N(0,\Sigma_t)}
\big[\hat L_{S_t}(W_{\mathrm{frozen}}+\hat{\Delta W}_t+\varepsilon)\big]
&=
\hat L_{S_t}(W_{\mathrm{frozen}}+\hat{\Delta W}_t) \\
&\quad +
\mathbb E_{\varepsilon \sim \mathcal N(0,\Sigma_t)}
\Big[
\hat L_{S_t}(W_{\mathrm{frozen}}+\hat{\Delta W}_t+\varepsilon)
-
\hat L_{S_t}(W_{\mathrm{frozen}}+\hat{\Delta W}_t)
\Big].
\end{align*}
The second term is exactly the expected Sharpness on the subspace
\[
\mathrm{E\text{-}Sh}_t^\Delta(\hat{\Delta W}_t,\Sigma_t)
:=
\mathbb E_{\varepsilon \sim \mathcal N(0,\Sigma_t)}
\Big[
\hat L_{S_t}(W_{\mathrm{frozen}}+\hat{\Delta W}_t+\varepsilon)
-
\hat L_{S_t}(W_{\mathrm{frozen}}+\hat{\Delta W}_t)
\Big],
\]
completes the proof.
\end{proof}



\subsection{Proof of Theorem \ref{thm:pecl-bound}}
\label{proof:theorem}

% \subsection{Proof of }

\begin{lemma}[Pathwise PAC--Bayes with time-varying hyper-posterior]
\label{the: full}
For any $\delta\in(0,1]$, with probability at least $1-\delta$ over $S_{1:T}$,
\begin{equation}
\begin{aligned} 
& \frac{1}{T} \sum_{t=1}^T \underbrace{\mathbb{E}_{P_t \sim \mathcal{P}_t} \mathbb{E}_{W \sim Q_t(\cdot \mid P_t, S_t)} L_{\mathcal{D}_t}(W)}_{\text{test risk at step } t} 
\le \ \frac{1}{T} \sum_{t=1}^T \underbrace{\mathbb{E}_{P_t \sim \mathcal{P}_t} \mathbb{E}_{W \sim Q_t(\cdot \mid P_t, S_t)} \hat{L}_{S_t}(W)}_{\text{empirical risk at step } t} \\ 
& + \sqrt{\frac{ \sum_{t=1}^T \mathbb{E}_{P_t \sim \mathcal{P}_t} \mathrm{KL}\big(Q_t(\cdot \mid P_t, S_t) \,\|\, P_t\big) + \sum_{t=1}^T \mathrm{KL}(\mathcal{P}_t \,\|\, \mathcal{P}_{t-1}) + \log\frac{1}{\delta} }{2 \sum_{t=1}^T (m_t - 1)}}. 
\end{aligned} 
\end{equation}
\end{lemma}


\begin{proof}
    


% \subsection*{A.1\quad Two-level variational inequalities}
Let $B_T:=\sum_{t=1}^T (m_t-1)$ and define $\alpha_t:=(m_t-1)/B_T$ so that $\sum_{t=1}^T\alpha_t=1$.
% \paragraph{Goal quantity.}
Define the step-$t$ generalization gap
\[
\Delta_t
:=\mathbb E_{P_t\sim\mathcal P_{t-1}}\,
\mathbb E_{W\sim Q_t(\cdot\mid P_t,S_t)}
\big(L_{\mathcal D_t}(W)-\hat L_{S_t}(W)\big),
\]
and the size-weighted gap $\bar\Delta:=\sum_{t=1}^T \alpha_t \Delta_t$.
Fix $\lambda>0$ and set
\(
g_t(W):=\lambda\alpha_t\big(L_{\mathcal D_t}(W)-\hat L_{S_t}(W)\big).
\)

\paragraph{Step 1 Model level.}
For any fixed $P$ and $S_t$,
\begin{equation}
\label{eq:dv-w}
\mathbb E_{W\sim Q_t(\cdot\mid P,S_t)} g_t(W)
\;\le\;
\mathrm{KL}\!\big(Q_t(\cdot\mid P,S_t)\,\|\,P\big)
+\log\mathbb E_{W\sim P}\!\big[e^{g_t(W)}\big].
\end{equation}
Taking $\mathbb E_{S_t}$ and using the standard Hoeffding--MGF bound for bounded losses,
\begin{equation}
\label{eq:mgf-mean-gap}
\mathbb E_{S_t}\log\mathbb E_{W\sim P} e^{g_t(W)}
\;\le\; \frac{(\lambda\alpha_t)^2}{8\,(m_t-1)}.
\end{equation}
Therefore,
\begin{equation}
\label{eq:model-step}
\mathbb E_{S_t}\mathbb E_{W\sim Q_t(\cdot\mid P,S_t)} g_t(W)
\;\le\;
\mathbb E_{S_t}\mathrm{KL}\!\big(Q_t(\cdot\mid P,S_t)\,\|\,P\big)
+\frac{(\lambda\alpha_t)^2}{8\,(m_t-1)}.
\end{equation}
 % (change of measure in $P$-space)
\paragraph{Step 2 Hyper level.} 
Define
\[
h_t(P):=\mathbb E_{S_t}\mathbb E_{W\sim Q_t(\cdot\mid P,S_t)} g_t(W)
-\mathbb E_{S_t}\mathrm{KL}\!\big(Q_t(\cdot\mid P,S_t)\,\|\,P\big).
\]
By \eqref{eq:model-step}, $h_t(P)\le \frac{(\lambda\alpha_t)^2}{8\,(m_t-1)}$ for all $P$.
Applying DV in $P$-space to change the measure from $\mathcal P_{t-1}$ to $\mathcal P_t$ gives
\begin{equation}
\label{eq:dv-p}
\log\mathbb E_{P\sim\mathcal P_t}\!\big[e^{h_t(P)}\big]
\;\le\; \mathrm{KL}(\mathcal P_t\|\mathcal P_{t-1})
+\log\mathbb E_{P\sim\mathcal P_{t-1}}\!\big[e^{h_t(P)}\big]
\;\le\; \mathrm{KL}(\mathcal P_t\|\mathcal P_{t-1})
+\frac{(\lambda\alpha_t)^2}{8\,(m_t-1)}.
\end{equation}
Using $\log\mathbb E[e^X]\ge \mathbb E[X]$ yields
\begin{equation}
\label{eq:p-expect}
\mathbb E_{P_t\sim\mathcal P_t} h_t(P_t)
\;\le\; \mathrm{KL}(\mathcal P_t\|\mathcal P_{t-1})
+\frac{(\lambda\alpha_t)^2}{8\,(m_t-1)}.
\end{equation}
Unfolding $h_t(P)$ and interchanging expectations (Fubini),
\begin{equation}
\label{eq:two-term}
\mathbb E_{S_t}\Big[
\lambda\alpha_t\Delta_t
-\alpha_t\,\mathbb E_{P_t\sim\mathcal P_{t-1}}\mathrm{KL}\!\big(Q_t(\cdot\mid P_t,S_t)\,\|\,P_t\big)
\Big]
\;\le\;
\mathrm{KL}(\mathcal P_t\|\mathcal P_{t-1})
+\frac{(\lambda\alpha_t)^2}{8\,(m_t-1)}.
\end{equation}

% \subsection*{A.2\quad Conditional MGF and supermartingale}
\paragraph{Step 3 Martingale}
From \eqref{eq:two-term} and the tower property, one obtains the conditional exponential bound
\begin{equation}
\label{eq:cond-mgf}
\mathbb E\!\left[
\exp\!\Big(
\lambda\alpha_t\Delta_t
-\alpha_t\,\mathbb E_{P_t\sim\mathcal P_{t-1}}\mathrm{KL}(Q_t\|P_t)
-\mathrm{KL}(\mathcal P_t\|\mathcal P_{t-1})
-\frac{(\lambda\alpha_t)^2}{8\,(m_t-1)}
\Big)\,\Big|\,\mathcal F_{t-1}\right]\;\le\;1.
\end{equation}
Define $X_t$ to be the exponential increment inside \eqref{eq:cond-mgf} and
\[
\mathcal{M}_0:=1,\qquad
\mathcal{M}_t:=\prod_{s=1}^t X_s.
\]
Then $(\mathcal{M}_t)_{t\ge0}$ is a nonnegative supermartingale, so
$\mathbb E[\mathcal{M}_T]\le 1$.
By Markov's inequality, with probability at least $1-\delta$,
\begin{equation}
\label{eq:hpp}
\sum_{t=1}^T \lambda\alpha_t\Delta_t
\;\le\;
\sum_{t=1}^T \alpha_t\,\mathbb E_{P_t\sim\mathcal P_{t-1}}\mathrm{KL}(Q_t\|P_t)
\;+\; \sum_{t=1}^T \mathrm{KL}(\mathcal P_t\|\mathcal P_{t-1})
\;+\; \frac{\lambda^2}{8}\sum_{t=1}^T \frac{\alpha_t^2}{(m_t-1)}
\;+\; \log\tfrac1\delta.
\end{equation}

% \subsection*{A.3\quad Size-weighted self-normalization and optimization}

By $\alpha_t=(m_t-1)/B_T$,
\[
\sum_{t=1}^T \frac{\alpha_t^2}{(m_t-1)}
=\sum_{t=1}^T \frac{(m_t-1)^2/B_T^2}{(m_t-1)}
=\frac{1}{B_T}.
\]
Hence, \eqref{eq:hpp} reads, for all $\lambda>0$ and with probability at least $1-\delta$,
\begin{equation}
\label{eq:gap-preopt}
\bar\Delta
\;\le\;
\frac{1}{\lambda}\Big(
\sum_{t=1}^T \alpha_t\,\mathbb E_{P_t\sim\mathcal P_{t-1}}\mathrm{KL}(Q_t\|P_t)
+ \sum_{t=1}^T \mathrm{KL}(\mathcal P_t\|\mathcal P_{t-1})
+ \log\tfrac1\delta\Big)
\;+\;
\frac{\lambda}{8B_T}.
\end{equation}
Optimizing the right-hand side over $\lambda>0$ (take $\lambda^\star=\sqrt{8B_T\,A}$ with
$A:=\sum_{t=1}^T \alpha_t\,\mathbb E_{P_t\sim\mathcal P_{t-1}}\mathrm{KL}(Q_t\|P_t)
+ \sum_{t=1}^T \mathrm{KL}(\mathcal P_t\|\mathcal P_{t-1}) + \log\tfrac1\delta$) yields
\begin{equation}
\label{eq:weighted-final}
\bar\Delta
\;\le\;
\sqrt{\frac{A}{2B_T}}
\;\le\;
\sqrt{\frac{
\sum_{t=1}^T \mathbb E_{P_t\sim\mathcal P_{t-1}}\mathrm{KL}(Q_t\|P_t)
+ \sum_{t=1}^T \mathrm{KL}(\mathcal P_t\|\mathcal P_{t-1})
+ \log\tfrac1\delta}{2\,\sum_{t=1}^T (m_t-1)}}.
\end{equation}
The last inequality uses $\sum_t \alpha_t x_t \le \sum_t x_t$ for nonnegative $(x_t)$.

\paragraph{Step 4 From weighted to the statement of Lemma~\ref{the: full}.}
The theorem controls the unweighted average $\frac1T\sum_{t=1}^T \Delta_t$ and uses the smoothed empirical risk on the right-hand side.
When $m_t$ are homogeneous, $\alpha_t\equiv 1/T$ and $\bar\Delta=\frac1T\sum_t\Delta_t$.
For heterogeneous $m_t$, we keep the same right-hand side and upper-bound the numerator $\sum_t \alpha_t \mathbb E_{P_t}\mathrm{KL}(Q_t\|P_t)$ by $\sum_t \mathbb E_{P_t}\mathrm{KL}(Q_t\|P_t)$, which yields exactly the denominator $2\sum_t(m_t-1)$ in the theorem.

% \subsection*{A.4\quad Gaussian smoothing (LoRA) corollary}

If $Q_t(\cdot\mid P_t,S_t)=\mathcal N(\hat W_t,\Sigma_t)$ (e.g., supported on an adapter subspace), then
\[
\mathbb E_{W\sim Q_t(\cdot\mid P_t,S_t)}\hat L_{S_t}(W)
=\mathbb E_{\varepsilon\sim\mathcal N(0,\Sigma_t)}\hat L_{S_t}(\hat W_t+\varepsilon),
\]
which converts the empirical term into the smoothed empirical loss (expected sharpness) used in the main text; substituting this identity into \eqref{eq:weighted-final} yields the empirical term in Lemma~\ref{the: full}.

\end{proof}



\begin{theorem}[Sharpness-aware PAC--Bayes bound for PECL]
% \label{thm:pecl-bound}

For any $\delta\in(0,1]$, with probability at least $1-\delta$ over $S_1,\ldots,S_T$,
\begin{equation}\label{eq:pecl-main}
\begin{aligned}
&\frac{1}{T}\sum_{t=1}^T
\mathbb{E}_{P_t\sim\mathcal P_{t-1}}\mathbb{E}_{W\sim Q_t}L_{\mathcal{D}_t}(W)
\ \le\
\ \frac{1}{T}\sum_{t=1}^T
\mathbb{E}_{P_t\sim\mathcal P_{t-1}}
\mathbb{E}_{\varepsilon \sim \mathcal{N}(0, \Sigma_t)}[\hat{L}_{S_t}(W_0 + \hat{\Delta W}_t + \varepsilon)
\\
&\quad +\sqrt{\frac{
\displaystyle \sum_{t=1}^T \mathbb{E}_{P_t^\Delta\sim\mathcal P_{t-1}}\!\Big[\mathrm{KL}\!\big(Q_t^\Delta\|P_t^\Delta\big)\Big]
\;+\;\displaystyle \sum_{t=1}^T \mathrm{KL}\!\big(\mathcal P_t\|\mathcal P_{t-1}\big)
\;+\;\log\!\frac{1}{\delta}
}{2\,\displaystyle\sum_{t=1}^T (m_t-1)}}.
\end{aligned}
\end{equation}
\end{theorem}



\begin{proof}
The bound follows by specializing Theorem \ref{the: full} to the PECL setting through two key decompositions:
By Lemma~\ref{lem:kl-decomp-pecl}, the full-space KL terms reduce to their adapter-space counterparts:
\[
\mathrm{KL}(Q_t^{\mathrm{full}} \| P_t^{\mathrm{full}}) = \mathrm{KL}(Q_t^\Delta \| P_t^\Delta),
\quad
\mathrm{KL}(\mathcal{P}_t^{\mathrm{full}} \| \mathcal{P}_{t-1}^{\mathrm{full}}) = \mathrm{KL}(\mathcal{P}_t^\Delta \| \mathcal{P}_{t-1}^\Delta).
\]

% \textbf{Step 2: Sharpness decomposition via Lemma 6.}  
Under the Gaussian posterior $Q_t^\Delta = \mathcal{N}(\Delta\hat{W}_t, \Sigma_t)$, Lemma \ref{lem:subspace-loss-pecl} gives:
\[
\mathbb{E}_{W \sim Q_t^{\mathrm{full}}}[\hat{L}_{S_t}(W)] 
= \mathrm{E}_{\varepsilon \sim \mathcal{N}(0, \Sigma_t)}{L}_{S_t}(W_{t-1}+\Delta\hat{W}_t + \varepsilon)
 = \hat{L}_{S_t}(W_{t-1} + \Delta\hat{W}_t) + \mathrm{E\text{-}Sh}_t^\Delta(\Delta\hat{W}_t, \Sigma_t),
\]
where the perturbations are confined to the adapter subspace.

% \textbf{Step 3: Substitution into general bound.}
Substituting these decompositions into Lemma \ref{the: full} yields the final PECL-specific bound, where all complexity and sharpness terms are now restricted to the trainable LoRA subspace $\mathcal{W}_\Delta$.
% Substituting these decompositions into Theorem 2 yields the final PECL-specific bound, where all complexity and sharpness terms are now restricted to the trainable LoRA subspace $\mathcal{W}_\Delta$.

\end{proof}


\section{Implement detail}


Unless stated otherwise, we use SGD as the base optimizer, train each session for $20$ epochs and SAM radius $\rho = 0.05$. LoRA based variants are trained with learning rate $0.01$ without momentum or weight decay. The same training configuration is used across SeqLoRA, IncLoRA, and OLoRA, so that differences in performance can be attributed to the geometry of the adapter subspace and the choice of flatness optimizer rather than to changes in training budget.



    
\section{Supplementary results}

\subsection{Detailed results on ImageNet-R}


\begin{table}[htpb]
\centering
\caption{
Class incremental performance and weight space flatness metrics on ImageNet-R ($T=20$, $R=16$). 
The table reports standard class incremental metrics  together with three sharpness indices and empirical Hessian and Fisher statistics in the weight space.
}
\label{table:q1_ci_with_weight_flatness_reorg}
\setlength{\tabcolsep}{4pt}
\resizebox{\linewidth}{!}{
\begin{tabular}{l l ccc ccccc c}
\toprule
\multirow{2}{*}{\textbf{Method}} & \multirow{2}{*}{\textbf{Opt.}} &
\multicolumn{3}{c}{\textbf{CL ($\uparrow$)}} &
\multicolumn{2}{c}{\textbf{Sharpness ($\downarrow$)}} &
\multicolumn{2}{c}{\textbf{Emp Hessian ($\downarrow$)}} &
\multicolumn{2}{c}{\textbf{Emp Fisher ($\downarrow$)}} \\
\cmidrule(lr){3-5}\cmidrule(lr){6-7}\cmidrule(lr){8-9}\cmidrule(lr){10-11}
 & & $\operatorname{ACC}_{20}^{\mathrm{cls}}$  &$\operatorname{AAA}^{\mathrm{cls}}$   & $\operatorname{BWT}^{\mathrm{cls}}$
 & Sh$^{(0)}(\rho)$  & Sh$^{(1)}(\rho)$
 & $\lambda_{\max}(H)$ & $\mathrm{tr}(H)$  
 & $\lambda_{\max}(F)$ & $\mathrm{tr}(F)$   \\
\midrule
\multirow{1}{*}{\textbf{PerTaskLP}}
& SGD   & $55.22^{\pm 0.31}$ & $59.36^{\pm 0.19}$ & $-8.35^{\pm 0.75}$ & 0.072   & 0.072 & 5.31    & 1182.49  & 8.47   & 1015.07  \\
\midrule
\multirow{2}{*}{\textbf{SeqFT}}
& SGD   &$50.72^{\pm 2.20}$ & $59.89^{\pm 1.32}$ & $-38.94^{\pm 1.87}$ & 0.582  & 0.517 & 4692.68 & 25621.40 & 538.79 & 29703.66 \\
& SAM   &$56.07^{\pm 2.60}$ & $66.47^{\pm 1.53}$ & $-33.04^{\pm 4.71}$ & 0.259  & 0.252 & 474.25  & 19747.59 & 75.43  & 5881.32  \\
 &GAM  &$59.78^{\pm 0.74}$ & $67.31^{\pm 1.18}$ & $-32.70^{\pm 0.26}$ & 0.177  & 0.176 & 877.97  & 18879.15 & 32.00  & 3273.08  \\
\midrule
\multirow{2}{*}{\textbf{SeqLoRA}}
& SGD    &$63.14^{\pm 2.17}$ & $69.07^{\pm 1.33}$ & $-11.91^{\pm 1.82}$ & 0.142  & 0.114 & 47.37 & 1502.55 & 35.51 & 1401.15 \\
& SAM    &$67.04^{\pm 0.71}$ & $71.42^{\pm 0.71}$ & $-9.26^{\pm 0.66}$ & 0.120  & 0.093 & 24.41   & 996.59   & 21.39  & 821.45   \\
 &GAM  &$73.31^{\pm 0.52}$ & $76.80^{\pm 0.56}$ & $-9.37^{\pm 0.61}$ & 0.060  &0.058  & 4.81    & 339.49   & 3.72   & 240.91   \\
\midrule
\multirow{2}{*}{\textbf{IncLoRA}}
& SGD    &$65.77^{\pm 2.84}$ & $71.29^{\pm 0.78}$ & $-4.47^{\pm 2.85}$ & 0.076    & 0.062 & 15.46 & 969.34  & 8.89  & 836.91 \\
& SAM    &$68.86^{\pm 0.92}$ & $72.59^{\pm 0.39}$ & $-3.51^{\pm 0.50}$ & 0.053   & 0.052 & 3.72    & 712.88   & 4.35   & 526.49   \\
 &GAM  &$72.86^{\pm 0.23}$ & $76.52^{\pm 0.48}$ & $-6.41^{\pm 0.48}$ & 0.035   & 0.035 & 2.69    & 281.22   & 2.09   & 192.25   \\
\midrule
\multirow{2}{*}{\textbf{OLoRA}}
& SGD    &$66.64^{\pm 0.44}$ & $71.41^{\pm 0.47}$ & $-3.52^{\pm 0.43}$ & 0.078   & 0.058 & 6.94    & 938.93   & 9.00   & 831.23   \\
& SAM    &$68.53^{\pm 0.78}$ & $72.75^{\pm 0.76}$ & $-3.16^{\pm 0.68}$ & 0.066   & 0.048 & 4.40    & 775.40   & 7.41   & 628.12   \\
 &GAM &$72.44^{\pm 0.36}$ & $76.05^{\pm 0.52}$ & $-7.12^{\pm 0.88}$ & 0.036   & 0.036 & 2.22    & 261.47   & 2.08   & 198.30\\
\midrule
\multirow{1}{*}{\textbf{l2p}}
& Adam   & $71.35^{\pm 0.05}$ & $74.84^{\pm 0.73}$ & $-6.34^{\pm 0.50}$  &-  &- &- &- &- &- \\
\midrule
\multirow{1}{*}{\textbf{dualprompt}}
& Adam   & $71.59^{\pm 0.68}$ & $74.80^{\pm 1.18}$ & $-3.38^{\pm 0.45}$ &-  &- &- &- &- &- \\
\bottomrule
\end{tabular}
}
\end{table}


\begin{table}[htpb]
\centering
\caption{
NME based performance and feature space diagnostics on ImageNet-R ($T=20$, $R=16$). 
The table reports NME accuracy metrics together with prototype drift, CKA similarity to the initial representation, and spectral properties of the empirical feature matrix.
}
\label{tab:q1_lp_feat_reorg}
\setlength{\tabcolsep}{4pt}
\renewcommand{\arraystretch}{0.95} % 可选：稍微压缩行距
\resizebox{0.9\linewidth}{!}{
\begin{tabular}{l l ccc  cc cc}
\toprule
\multirow{2}{*}{\textbf{Method}} & \multirow{2}{*}{\textbf{Opt.}} &
\multicolumn{3}{c}{\textbf{NME} ($\uparrow$)} &
\multicolumn{2}{c}{\textbf{Feature Stability}} &
\multicolumn{2}{c}{\textbf{Feature Matrix}} \\
\cmidrule(lr){3-5} \cmidrule(lr){6-7} \cmidrule(lr){8-9}
&  & $ \operatorname{ACC}_{20}^{\mathrm{nme}}$  & $\operatorname{AAA}^{\mathrm{nme}}$  & $\operatorname{BWT}^{\mathrm{nme}}$ 
 & L2 drift & CKA$_{0}$
 & $\mathrm{tr}(E)$ & eff \\ 
\midrule
\multirow{1}{*}{\textbf{PerTaskLP}}
  & SGD &$62.96^{\pm0.12}$ & $66.58^{\pm0.37}$ & $-7.66^{\pm1.27}$
     & 0        & 1        & 0.82 & 128.11  \\
\midrule
\multirow{2}{*}{\textbf{SeqFT}}
 & SGD     &$69.49^{\pm0.16}$ & $75.10^{\pm0.57}$ & $-16.11^{\pm0.23}$ & 29.57 & 0.78 & 0.58 & 93.53  \\
 & SAM     &$72.36^{\pm2.46}$ & $78.31^{\pm0.50}$ & $-14.04^{\pm3.30}$ & 17.04 & 0.86 & 0.74 & 127.28 \\
 &GAM  &$73.74^{\pm0.88}$ & $78.32^{\pm0.94}$ & $-13.36^{\pm0.20}$ & 13.80 & 0.85 & 0.78 & 127.87 \\
\midrule
\multirow{2}{*}{\textbf{SeqLoRA}}
 & SGD     &$72.82^{\pm1.27}$ & $76.62^{\pm0.68}$ & $-6.75^{\pm1.29}$  & 16.58 & 0.73 & 0.77 & 125.50 \\
 & SAM     &$74.53^{\pm0.60}$ & $77.62^{\pm0.43}$ & $-5.54^{\pm0.72}$  & 14.82 & 0.73 & 0.81 & 135.38 \\
 &GAM   &$75.74^{\pm0.15}$ & $78.86^{\pm0.53}$ & $-6.19^{\pm0.33}$  & 11.28 & 0.78 & 0.78 & 133.96 \\
\midrule
\multirow{2}{*}{\textbf{IncLoRA}}
 & SGD      &$75.41^{\pm0.26}$ & $78.37^{\pm0.26}$ & $-4.32^{\pm0.59}$   & 16.87 & 0.69 & 0.67 & 121.18 \\
 & SAM     &$75.95^{\pm0.07}$ & $78.68^{\pm0.25}$ & $-3.87^{\pm0.03}$ & 15.14 & 0.68 & 0.70 & 127.16 \\
&GAM  &$76.72^{\pm0.54}$ & $79.42^{\pm0.40}$ & $-3.63^{\pm0.67}$  & 12.15 & 0.75 & 0.73 & 119.68 \\
\midrule
\multirow{2}{*}{\textbf{OLoRA}}
 & SGD     &$75.10^{\pm0.41}$ & $78.65^{\pm0.67}$ & $-3.90^{\pm0.54}$ & 17.23 & 0.69 & 0.68 & 120.89 \\
 & SAM     &$75.73^{\pm0.32}$ & $78.70^{\pm0.67}$ & $-3.47^{\pm0.54}$ & 15.17 & 0.72 & 0.70 & 123.76 \\
 &GAM  &$76.56^{\pm0.53}$ & $79.16^{\pm0.43}$ & $-3.37^{\pm0.68}$  & 12.16 & 0.75 & 0.74 & 119.36\\
\bottomrule
\end{tabular}
}
\end{table}

\begin{figure}[thbp]
    \centering
    \begin{minipage}{0.495\textwidth}
        \centering
        \includegraphics[width=0.9\linewidth]{figures_TRML/AAA_curve_Imagenet-R_summary_inc10_rank16—point2.pdf}
        % {figures_TRML/AAA_curve_Imagenet-R_summary_inc10_rank16_replot2.pdf}
        % \subcaption{Accuracy$\operatorname{ACC}_t^{\mathrm{cls}}$ on ImageNet-R}
        \label{fig:curves_imageR_1_AAA}
    \end{minipage}
    \begin{minipage}{0.495\textwidth}
        \centering
        \includegraphics[width=0.9\linewidth]{figures_TRML/NME_curve_Imagenet-R_summary_inc10_rank16_point2.pdf}
        % {figures_TRML/NME_curve_Imagenet-R_summary_inc10_rank16_replot2.pdf}
        % \subcaption{NME-based accuracy $\operatorname{ACC}_t^{\mathrm{nme}}$ on ImageNet-R}
        \label{fig:curves_imageR_2_NME}
    \end{minipage}
    \caption{
    % Summary curves of accuracy with the classifier head $\operatorname{ACC}_t^{\mathrm{cls}}$ and prototype-based accuracy $\operatorname{ACC}_t^{\mathrm{nme}}$ on ImageNet-R ($T=20, r=16$). 
    %     The curves report performance after each task for different class incremental methods and optimization configurations. 
        Summary curves of classifier-head accuracy $\operatorname{ACC}_t^{\mathrm{cls}}$ and prototype-based accuracy $\operatorname{ACC}_t^{\mathrm{nme}}$ on ImageNet-R ($T=20, r=16$) over the task sequence, comparing different class-incremental methods and optimization configurations.
        Flatness-promoting optimizers such as SAM and GAM consistently outperform SGD, and their advantage becomes more pronounced on later tasks, which suggests more stable representations and improved performance.
        % Flatness promotion optimizers (SAM and GAM) consistently outperform SGD, and their advantage tends to increase over later tasks, indicating more stable representations and better performance.
    }
    \label{fig:combined_curves_imageR}
\end{figure}




%%%%%%%%%%%%%%%%%%%%%%%%%%%%%%%%%%%%%%%%%%%%%%%%%%%%%%%%%%%%%%%%%%%%%%%%%%%%%%%
%%%%%%%%%%%%%%%%%%%%%%%%%%%%%%%%%%%%%%%%%%%%%%%%%%%%%%%%%%%%%%%%%%%%%%%%%%%%%%%



\textbf{EFM analyse}

\begin{figure}[htbp]
\centering
\begin{minipage}{0.9\linewidth}
\centering
\includegraphics[width=0.9\linewidth]{figures_TRML/finetune_42_gam-sam-sgd_efm_spectrum_compare.pdf}
% \subcaption{Eigenvalue Spectrum in Sequential Fine-tuning (SeqFT)}
\label{fig:EFM_seqFT}
\end{minipage}
\caption{
Comparison of EFM eigenvalue spectra in the feature space across incremental learning methods. The horizontal axis represents eigenvalue rank (sorted by magnitude), while the vertical axis shows eigenvalue magnitude. Only the top 200 eigenvalues are non-zero, with remaining dimensions at zero. SAM consistently produces flatter, heavier-tailed spectra compared to SGD across both configurations, indicating reduced energy concentration in the leading principal directions. This spectral profile suggests that SAM learns more isotropic representations with higher effective dimensionality.
}
\label{fig:EFM_spectra_comparison}
\end{figure}



\subsection{Detailed Resuts on ImageNet-C and ImageNet-P}


\begin{figure*}[tbp]
    \centering
    % ImageNet-C: cls / nme
    \begin{minipage}{0.495\linewidth}
        \centering
        \includegraphics[width=0.85\linewidth]{figures_TRML/AAA_curve_Imagenet-C_summary_inc10_rank16_ponit2.pdf}
        % \subcaption{ImageNet-C: $\operatorname{ACC}_t^{\mathrm{cls}}$}
        \label{fig:aaa_imagenet_c}
    \end{minipage}\hfill
    \begin{minipage}{0.495\linewidth}
        \centering
        \includegraphics[width=0.85\linewidth]{figures_TRML/NME_curve_Imagenet-C_summary_inc10_rank16_point2.pdf}
        % \subcaption{ImageNet-C: $\operatorname{ACC}_t^{\mathrm{nme}}$}
        \label{fig:nme_imagenet_c}
    \end{minipage}\hfill
    % ImageNet-P: cls / nme
    \begin{minipage}{0.495\linewidth}
        \centering
        \includegraphics[width=0.85\linewidth]{figures_TRML/AAA_curve_Imagenet-P_summary_inc10_rank16_point2.pdf}
        % \subcaption{ImageNet-P: $\operatorname{ACC}_t^{\mathrm{cls}}$}
        \label{fig:aaa_imagenet_p}
    \end{minipage}\hfill
    \begin{minipage}{0.495\linewidth}
        \centering
        \includegraphics[width=0.85\linewidth]{figures_TRML/NME_curve_Imagenet-P_summary_inc10_rank16_point2.pdf}
        % \subcaption{ImageNet-P: $\operatorname{ACC}_t^{\mathrm{nme}}$}
        \label{fig:nme_imagenet_p}
    \end{minipage}
    \caption{ 
        The full trajectories
of$\operatorname{ACC}_t^{\mathrm{nme}}$ on ImageNet-C and ImageNet-P. 
        % The four panels report performance after each task for the same class incremental methods and optimization configurations as in Figure~\ref{fig:combined_curves_imageR}. 
         SAM and GAM consistently outperform SGD on both benchmark, which indicates that promoting flat minima in the adapter improves robustness of the learned representations.
         %   and $\operatorname{ACC}_t^{\mathrm{cls}}$
    }
    % \label{fig:robustness_curves_comparison}
\end{figure*}

% \subsection{Ablation Study}


\subsection{Time consusme analyse}

\begin{table}[thpb]
\centering
\caption{Average per‑task training time (seconds) across methods and optimizers on ImageNet‑R}
\label{tab:lora_opt_metric}
\setlength{\tabcolsep}{4pt}
\renewcommand{\arraystretch}{1}
\resizebox{\linewidth}{!}{
\begin{tabular}{l ccccc}
\toprule
\textbf{Method} & \textbf{SGD} & \textbf{RWP} & \textbf{SAM} & \textbf{GAM} & \textbf{C-FLAT} \\
\midrule
SeqLoRA & 282.72 ± 61.34  & 244.95 ± 6.38  & 472.29 ± 83.70  & 670.20 ± 217.92 & 532.62 ± 35.93   \\
IncLoRA & 310.54 ± 141.00 & 658.34 ± 83.90 & 777.91 ± 239.29 & 1247.56 ± 5.05  & 1189.40 ± 109.93 \\
OLoRA   & 412.99 ± 146.84 & 644.19 ± 55.05 & 628.44 ± 181.29 & 836.99 ± 327.65 & 990.88 ± 360.79  \\
\bottomrule
\end{tabular}
}
\end{table}






\end{document}


% This document was modified from the file originally made available by
% Pat Langley and Andrea Danyluk for ICML-2K. This version was created
% by Iain Murray in 2018, and modified by Alexandre Bouchard in
% 2019 and 2021 and by Csaba Szepesvari, Gang Niu and Sivan Sabato in 2022.
% Modified again in 2023 and 2024 by Sivan Sabato and Jonathan Scarlett.
% Previous contributors include Dan Roy, Lise Getoor and Tobias
% Scheffer, which was slightly modified from the 2010 version by
% Thorsten Joachims & Johannes Fuernkranz, slightly modified from the
% 2009 version by Kiri Wagstaff and Sam Roweis's 2008 version, which is
% slightly modified from Prasad Tadepalli's 2007 version which is a
% lightly changed version of the previous year's version by Andrew
% Moore, which was in turn edited from those of Kristian Kersting and
% Codrina Lauth. Alex Smola contributed to the algorithmic style files.
